\documentclass[11pt]{report}

\usepackage[T1]{fontenc}
\usepackage[utf8]{inputenc}
\usepackage[a4paper,margin=40mm]{geometry}

\usepackage[english]{babel}
\usepackage{microtype}
\usepackage{sourcecodepro}
\usepackage{indentfirst}
\usepackage{amsmath}
\usepackage{amsthm}
\usepackage{amssymb}
\usepackage{mathtools}
\usepackage{xcolor}
\usepackage{multirow}
\usepackage{hyperref}

\usepackage[
	font={footnotesize,it},
	justification=raggedright,
	singlelinecheck=true,
	labelfont={up,sc,bf},
	labelsep=colon,
	position=auto,
	skip=0pt
]{caption}

\usepackage{titlesec}
\titleformat{\part}[display]
	{\scshape\Huge\filcenter}
	{\partname~\thepart}
	{0pt}
	{}
\titleformat{\chapter}[display]
	{\scshape\Huge\filcenter}
	{\chaptertitlename~\thechapter}
	{0pt}
	{}
\titleformat{\section}
	{\scshape\huge}
	{\llap{\thesection~}}
	{0pt}
	{}
\titleformat{\subsection}
	{\scshape\LARGE}
	{\llap{\thesubsection~}}
	{0pt}
	{}
\titleformat{\subsubsection}
	{\scshape\Large}
	{\llap{\thesubsubsection~}}
	{0pt}
	{}
\titleformat{\paragraph}[runin]
	{\bfseries\scshape\large}
	{\llap{\theparagraph~}}
	{0pt}
	{}
\titleformat{\subparagraph}[runin]
	{\bfseries\scshape\normalsize}
	{\llap{\thesubparagraph~}}
	{0pt}
	{}
\titlespacing*{\chapter}{0pt}{0pt}{2.0\baselineskip}
\titlespacing*{\paragraph}{\parindent}{3.25ex plus 1ex minus .2ex}{1em}
\titlespacing*{\subparagraph}{2\parindent}{3.25ex plus 1ex minus .2ex}{1em}

\usepackage{fancyhdr}
\renewcommand{\chaptermark}[1]{\markboth{\thechapter~#1}{}}
\renewcommand{\sectionmark}[1]{\markright{\thesection~#1}}
\fancyhead{}
\fancyhead[L]{\nouppercase{\scshape\leftmark~//~\rightmark}}
% \fancyhead[R]{\nouppercase{\scshape\rightmark}}
\fancyfoot{}
\fancyfoot[C]{\thepage}
\renewcommand{\footrulewidth}{0.0pt}
\renewcommand{\headrulewidth}{0.5pt}
\fancyheadoffset{2cm}
\fancyfootoffset{2cm}
\setlength{\headheight}{14pt}

% \AddToHook{cmd/section/before}{\ifnum\value{section}>0 \clearpage\fi}

\setcounter{secnumdepth}{3}
\setcounter{tocdepth}{3}

\usepackage[newfloat=true]{minted}
\setminted{
	fontsize=\scriptsize,
	bgcolor=gray!10,
	tabsize=4,
	tabcolor=gray!50,
	spacecolor=gray!50,
	autogobble=true,
	ignorelexererrors=true,
	numbersep=4pt,
	linenos=true,
	bgcolorpadding=2pt,
	% style=autumn,
	% style=sas,
	% style=murphy,
	% style=tango,
	% style=friendly_grayscale
	style=algol_nu,
	% style=bw,
}
\setminted[Bash]{
	linenos=false,
}
\setminted[antlr]{
	linenos=false,
}
\AddToHook{env/listing/begin}{\setminted{vspace=0pt}}
\newcommand{\giovilexer}{./giovi_pygments/giovi.py:GioviLexer}

\usepackage{float}
\floatplacement{listing}{H}

\usepackage{marginnote}
\newcommand{\comment}[1]{\marginnote{\color{black!30}\small\sffamily#1}}
\newcommand{\commentinline}[1]{\parbox{0.9\linewidth}{\color{black!30}\small\sffamily#1}}
\setlength{\marginparwidth}{3cm}

\newcommand{\See}[1]{\S\ref{#1}~\textsc{\nameref{#1}}}
\newcommand{\see}[1]{\S\ref{#1}}
\newcommand{\set}[1]{\left\{#1\right\}}
\newcommand{\glist}[1]{\left[#1\right]}
\newcommand{\map}[1]{\left[#1\right]}
\newcommand{\infrule}[3][gathered]{\frac{\begin{#1}#2\end{#1}}{\begin{gathered}#3\end{gathered}}}
\newcommand{\oprule}[3][gathered]{\frac{\begin{#1}#2\end{#1}}{\begin{gathered}#3\end{gathered}}}
% \newcommand{\infrule}[2]{\colorbox{gray!2}{$\frac{\begin{gathered}#1\end{gathered}}{\begin{gathered}#2\end{gathered}}$}}
\newcommand{\semi}{;\ }
\newcommand{\semiq}{;\quad}
\newcommand{\then}{\vdash}
\newcommand{\Then}{\vDash}
\newcommand{\tuple}[1]{\left<#1\right>}
\newcommand{\code}[1]{\colorbox{gray!10}{\vrule height 7pt depth 0pt width 0pt \texttt{\scriptsize#1}}}
\newcommand{\nb}[1]{\textbf{N.B. } #1}
\newcommand{\narreq}{\leq:}
\newcommand{\eval}{\equiv}
\newcommand{\bop}{\diamond}
\newcommand{\alloc}{\mathrm{alloc}}
% \newcommand{\copi}{\mathrm{copy}}
\newcommand{\addr}{\alpha}
\newcommand{\expr}{\epsilon}
\newcommand{\val}{\nu} % eventually use chi
\newcommand{\pattern}{p} % eventually use pi
\newcommand{\Id}{\mathcal{I}}
\newcommand{\Expr}{\mathcal{E}}
\newcommand{\Val}{\mathcal{V}}
\newcommand{\Zero}{\mathcal{Z}}
\newcommand{\Pattern}{\mathcal{P}}
\newcommand{\PatternTerminal}{\Pattern_T}
\newcommand{\PatternProduct}{\Pattern_P}
\newcommand{\PatternIterator}{\Pattern_I}
\newcommand{\decons}{\delta}
\newcommand{\self}{\texttt{self}}
\newcommand{\Self}{\texttt{Self}}
\newcommand{\Number}{\texttt{Number}}
\newcommand{\Product}{\texttt{Product}}
\newcommand{\Type}{\texttt{Type}}
\newcommand{\Bool}{\texttt{Bool}}
\newcommand{\Text}{\texttt{Text}}
\newcommand{\Void}{\texttt{Void}}
\newcommand{\Int}{\texttt{Int}}
\newcommand{\Natural}{\texttt{Natural}}
\newcommand{\Real}{\texttt{Real}}
\newcommand{\List}{\texttt{List}}
\newcommand{\Map}{\texttt{Map}}
\newcommand{\Function}{\texttt{Function}}
\newcommand{\Tuple}{\texttt{Tuple}}
\newcommand{\Record}{\texttt{Record}}
\newcommand{\Struct}{\texttt{Struct}}
\newcommand{\Range}{\texttt{Range}}
% \newcommand{\Module}{\texttt{Module}}
\newcommand{\env}[1]{[\![#1]\!]}
\newcommand{\Env}{\Xi}
\newcommand{\EnvSym}{\Gamma}
\newcommand{\EnvVar}{\Upsilon}
\newcommand{\EnvConst}{\Pi}
\newcommand{\EnvFunc}{\Phi}
\newcommand{\EnvStucts}{\Sigma}
% \newcommand{\EnvModules}{M}
\newcommand{\EnvLoc}{\Lambda}
\newcommand{\EnvStore}{\Psi}

\theoremstyle{definition}
\newtheorem{definition}{Definition}[section]

\makeatletter
\providecommand{\bigsqcap}{\mathop{\mathpalette\@updown\bigsqcup}}
\newcommand*{\@updown}[2]{\rotatebox[origin=c]{180}{$\m@th#1#2$}}
\makeatother

\renewcommand{\arraystretch}{1.5} % increase space between table rows

\begin{document}

% \setlength{\parindent}{0pt}
\setlength{\doublerulesep}{0.40cm}
% \setlength{\baselineskip}{11pt}
\abovedisplayshortskip=8pt
\belowdisplayshortskip=8pt

\begin{titlepage}
	\centering
	\includegraphics[width=5cm]{GioviAssets/logo/giovi_cropped.png}

	\vspace{1cm}
	{\Huge\bfseries Giovi Language Specification}

	\vspace{1cm}
	{\huge Simone Vicinanza}

	\vfill
	{\Large 2069-69-69}

\end{titlepage}

\tableofcontents

\clearpage

\pagestyle{fancy}

\part{The language}

\chapter{Introduction}

\section{What is Giovi}

Giovi (pronounced as \textit{"jovi"}) is a high level, multipurpose,
strongly typed, imperative programming and scripting language.

\begin{listing}
	\label{lst:hello_world}
	\begin{minted}{\giovilexer}
"hello world" -> _out
	\end{minted}
	\caption{The usual sample program that prints "hello world" to stdout}
\end{listing}

Giovi is originally born as a configuration language and as such, it relies on
a mostly structural type system; it puts a lot of emphasis on the use of few
basic data structures that, combined with simple yet powerful syntactic
features, allows writing code that is easy to read and fast to write.

\begin{listing}
	\label{lst:intro_overview}
	\caption[Syntax intro overview]{
		A variable declaration of type \code{User}, initialized with a variety
		of data types: primitive types (numbers, strings, boolean), collection
		types (lists, maps), product types and functions
	}
	\begin{minted}{\giovilexer}
var simonvic <- User.<
	id: 69
	username: "simonvic"
	aliaes: ["Simone", "Simon"]
	bio: "{
		hello I'm simonvic,
		and I'm working on Giovi!
		}"
	age: 26
	favoriteColor: Rgb.<0, 255, 0>
	credits: 69.420
	roles: [
		"admin" => true
		"devel" => true
		"guest" => false
	]
	onLogin: function(user User) {
		"hello ${user.username}" -> _out
		if "admin" in user.roles do {
			"Granting admin priviledges" -> _out
		}
	}
>
	\end{minted}
\end{listing}

Giovi evolved to be more than a configuration language, by providing
user-defined data structures, first-class functions and types. Giovi wants to
experiment with the idea of non-nullability, where everything has a
(overwritable) "default" value that, along the support for automatic data
widening (but explicit narrowing), allows writing flexible code with few
keystrokes.

\begin{listing}
	\label{lst:calculator}
	\caption{
		An example of a \code{calculator} function and the use of functions as
		first-class citizens
	}
	\begin{minted}{\giovilexer}
plus(x Real, y Real) => Real { return x + y }
minus(x Real, y Real) => Real { return x + y }
times(x Real, y Real) => Real { return x * y }
divided(x Real, y Real) => Real { return x / y }

#? Calculator function
#? leftHandSide: first operand
#? strategy: a function with signature (Real,Real)=>Real
#? rightHandSide: second operand
calculator(
	leftHandSide Real,
	strategy function(x Real, y Real) => Real,
	rightHandSide Real,
) => Real {
	# apply strategy
	return strategy(leftHandSide, rightHandSide)
}

calulator(5, plus, 2) -> _out # 5 + 2
calulator(5, minus, 2) -> _out # 5 - 2
calulator(5, times, 2) -> _out # 5 * 2
calulator(5, divided, 2) -> _out # 5 / 2
	\end{minted}
\end{listing}

The grammar design philosophy is to have a language that is easy to read and
fast to write. The syntax aims to have minimal lexical and syntactic clutter
where keywords are preferred over symbols, and to provide many quality-of-life
features that speed up the writing process, without hindering the reading by
relying on the use of symbols that is consistent and predictable.

\section{Executing Giovi code}

Giovi code can be executed from a file

\begin{minted}{Bash}
giovi some_script.giovi
\end{minted}

Alternatively, the source code can be passed directly to the CLI

\begin{minted}{Bash}
giovi -f '"hello world" -> _out'
\end{minted}

\subsection{REPL}

Giovi also has a minimal REPL (Read Evaluation Print Loop), that allows you to
write and execute code on the fly

\begin{minted}{Bash}
giovi --repl
> "The sum of 10 and 20 is equal to ${10 + 20}" -> _out
> The sum of 10 and 20 is equal to 30
\end{minted}

\chapter{Grammar}

\section{Introduction}

This chapter describes the lexicon and syntax of the language. Grammar rules
will be shown using a dialect of eBNF.

\begin{listing}
	\label{lst:ebnf_dialect}
	\caption{Example of a definition of \code{someRule} using the eBNF dialect}
	\begin{minted}{antlr}
someRule
	: 'some_token' someLabel=someOtherRule
	| zeroOrOneOfThis? zeroOrMoreThis* oneOrMoreThis+
	| ('some_token' | 'some_other_token')
	;
	\end{minted}
\end{listing}

\nb{
	Only a subsest of the rules will be shown, and they may not exactly match
	the real grammar. For the actual rules consult the appropriate grammar
	files.
}

\paragraph{Whitespaces} Whitespaces (e.g. spaces, tabs, newlines) are
determinant \textbf{only} during the lexical parsing, meaning that \code{!=} is
different than \code{! =}, where the former gets parsed as a single token while
the latter gets parsed as two different tokens; in any other case, whitepsaces
are ignored.

\section{Lexicon}

\subsection{Keywords}

\code{if}, \code{do}, \code{else}, \code{when}, \code{for}, \code{in},
\code{repeat}, \code{break}, \code{skip}, \code{return}\\*
\code{default}, \code{import}, \code{module}\\*
\code{const}, \code{var}, \code{alias},
\code{self}, \code{function}, \code{tuple}, \code{record}, \code{struct},
\code{type}, \code{and}, \code{or}, \code{not}, \code{as}

\paragraph{Reserved keywords} Some keywords are reserved for future-proofing
to ensure backward compatibility in the eventuality they get used.

\code{def}, \code{let}, \code{ref}, \code{at}, \code{with}, \code{interface},
\code{enum}

\subsection{Operators}

\begin{table}[H]
\resizebox{\linewidth}{!}{
	\begin{tabular}{rcl}
		\textbf{Operator}               & \textbf{Token}  & \textbf{Description}                            \\\hline
		PLUS                            & \texttt{+}      & Arithmetic sum, concatenation                   \\\hline
		MINUS                           & \texttt{-}      & Arithmetic subtraction, intersection            \\\hline
		MULTIPLICATION                  & \texttt{*}      & Arithmetic multiplication                       \\\hline
		DIVISION                        & \texttt{/}      & Arithmetic division                             \\\hline
		MODULO                          & \texttt{\%}     & Arithmetic remainder                            \\\hline
		EXPONENTIAL                     & \texttt{\^}     & Arithmetic power, text repetition               \\\hline
		EQUAL                           & \texttt{=}      & Loose equality                                  \\\hline
		NOT EQUAL                       & \texttt{!=}     & Loose inequality                                \\\hline
		LESS EQUAL                      & \texttt{<=}     & Arithmetic and lexicographical comparison       \\\hline
		LESS                            & \texttt{<}      & Arithmetic and lexicographical comparison       \\\hline
		GREATER                         & \texttt{>}      & Arithmetic and lexicographical comparison       \\\hline
		GREATER EQUAL                   & \texttt{>=}     & Arithmetic and lexicographical comparison       \\\hline
		CONTAINS                        & \texttt{=\~{}}  & Checks if a set of elements contains an element \\\hline
		IS CONTAINED                    & \texttt{\~{}=}  & Reverse of CONTAINS                             \\\hline
		\multirow{2}{*}{BOOL NOT}       & \texttt{not}    & \multirow{2}{*}{Boolean logic negation}         \\
	                                	& \texttt{!}      &                                                 \\\hline
		\multirow{2}{*}{BOOL AND}       & \texttt{and}    & \multirow{2}{*}{Boolean logic and}              \\
	                                	& \texttt{\&}     &                                                 \\\hline
		\multirow{2}{*}{BOOL OR}        & \texttt{or}     & \multirow{2}{*}{Boolean logic or}               \\
	                                	& \texttt{|}      &                                                 \\\hline
		BITWISE NOT                     & \texttt{!!}     & Bitwise negation                                \\\hline
		BITWISE AND                     & \texttt{\&\&}   & Bitwise and                                     \\\hline
		BITWISE OR                      & \texttt{||}     & Bitwise or                                      \\\hline
		BITWISE XOR                     & \texttt{|||}    & Bitwise exclusive or                            \\\hline
		BITWISE LEFT SHIFT              & \texttt{\{<\}}  & Bitwise left shift                              \\\hline
		BITWISE RIGHT SHIFT             & \texttt{\{>\}}  & Bitwise right shift                             \\\hline
		BITWISE RIGHT UNSIGNED SHIFT    & \texttt{\{>>\}} & Bitwise right unsigned shift                    \\\hline
		IS                              & \texttt{is}     & Runtime type check, pattern matching            \\\hline
		AS                              & \texttt{as}     & Explicit type coercion                          \\\hline
	    \\
		\multicolumn{3}{c}{Pseudo-Operators}                                                                \\\hline
		DOT ACCESS                      & \texttt{.}      & Members access                                  \\\hline
		DOTREF ACCESS                   & \texttt{::}     & Function member access                          \\\hline
	\end{tabular}
}
\end{table}

\subsection{Other tokens}

\code{..}, \code{:}, \code{,}, \code{;}, \code{\#}, \code{(}, \code{)}, \code{[}, \code{]},
\code{\{}, \code{\}}, \code{<}, \code{>}

\subsection{Comments}

There are two kinds of comments: single line and multiline.

\subsubsection{Single line comment}

\begin{minted}{antlr}
	comment: '#' comment_content? ('\n' | EOF) ;
	comment_content: [^\n]+ ;
\end{minted}

Every character after the comment token \code{\#} until a newline \code{\\n} is
considered part of the comment.

\subsubsection{Multiline comment}

\begin{minted}{antlr}
	comment_block: '#{' .* ('}#' | EOF) ;
\end{minted}

Every character between and the open comment token \code{\#\{} and the close
comment token \code{\}\#} is considered part of the comment.\\*
A multiline comment can also be inlined in the code.

\begin{minted}{\giovilexer}
if true #{ this is a comment }# do something()
\end{minted}

\paragraph{Conventions} By convention, \code{\#} is immediately followed by
a space, as the first character is reserved for "special" comments, for example
documentation: \code{\#? This is a documentation comment}

\comment{TODO: Specificare RFC riguardo language server e altro?}

\clearpage

\subsection{Literals}

\subsubsection{Boolean}

\begin{minted}{antlr}
	literal_boolean: 'true' | 'false' ;
\end{minted}

A \code{Bool} literal is simply expressed by \code{true} or \code{false}

\subsubsection{Text}
\label{sec:text_syntax}

\begin{minted}{antlr}
	literal_normal_text
		: '"' ([^\n] | escape_sequence | interpolation)* '"'
		;

	literal_normal_text_block
		: '"{' (. | escape_sequence | interpolation)* '}"'
		;

	literal_exact_text
		: "'" ([^\n] | escape_sequence)* "'"
		;

	literal_exact_text_block
		: "'{" (. | escape_sequence)* "}'"
		;
\end{minted}

There are four types of \code{Text} literals (a.k.a. strings).

\begin{itemize}
	\item \textbf{Normal text}: delimited by double quotes \code{"}
	\item \textbf{Exact text}: delimited by single quptes \code{'}
	\item \textbf{Normal text block}: delimited by \code{"\{} and \code{\}"}
	\item \textbf{Exact text block}: delimited by \code{'\{} and \code{\}'}
\end{itemize}

\paragraph{Differences} All types of text literals can contain escape
sequences.\\*
Normal and exact text must be terminated on the same line (they cannot contain
a newline).\\*
Normal and exact text blocks can span over multiple lines and do not require
the escape of the text delimiter.\\*
Normal text and normal text block allow interpolation.

\paragraph{Escape sequences} Some special characters may not appear in the
text, unless when escaped by prefixing it with \code{$\backslash$} (backslash).

\begin{minted}{antlr}
	escape_sequence: character_escape | unicode_escape ;
	character_escape: '\' [btnfrs"'\\${}] ;
	unicode_escape: '\' 'u' HexDigit HexDigit HexDigit HexDigit ;
\end{minted}

\paragraph{Interpolation} Normal text and normal text blocks support
interpolation, allowing the "injecttion" of a variable or expression in the
text.\\*
To inject a variable, simply prefix it with \code{\$} (dollar):\\*
\code{"My name is \$name"}.\\*
In order for an expression to be injected, it needs to be wrapped in braces:\\*
\code{"My name is \$\{user.getName()\}"}

\paragraph{Text blocks} Normal and exact text blocks allow the text to span
over multiple lines, without the need to escape the text delimiter.

\begin{listing}
	\label{lst:textblock_example}
	\begin{minted}{\giovilexer}
if true do {
	if true do {
		html Text <- "{
			<html>
				<head></head>
				<body>
					<p class="foo">Hello world</p>
				<body>
			</html>
			}"
	}
}
	\end{minted}
	\caption{Example of a textblock appearing indented in a nested scope}
\end{listing}

The last newline character will be included. If you don't want it, put the
closing tag on the previous line

\begin{minted}{\giovilexer}
t Text <- "{
	hello world}"
\end{minted}

\clearpage

\subparagraph{Incidental whitespaces} In a text block, white spaces (spaces,
newlines, tabs) are preserved, unless they're considered "incidental".\\*
Incidental white spaces are those needed for code indentation.

\begin{minted}[showtabs,tab=>,escapeinside=||]{\giovilexer}
if true do {
	if true do {
		html Text <- "{
			<html>
				<head></head>
				<body>
					<p>Hello world</p>
				<body>
			</html>
			}"
#^^^^^^^^^^^
# the first three tabulations are considered "incidental"
# and won't be included in the text
	}
}
\end{minted}

The amount of "incidental" whitespaces is defined by the line that has less of
them.

\nb{
	Only the highlighted part will actually be included in the text literal
}

\begin{listing}
	\label{lst:textblock_incidental_whitespaces}
	\begin{minted}[autogobble=false,showtabs,tab=>,escapeinside=||]{text}
			_out <- "{
			|\colorbox{gray!80}{<html>}|
			|\colorbox{gray!80}{	<head></head>}|
			|\colorbox{gray!80}{	<body>}|
			|\colorbox{gray!80}{		<p>Hello world</p>}|
			|\colorbox{gray!80}{	<body>}|
			|\colorbox{gray!80}{</html>}|
			}"

			_out <- "{
		|\colorbox{gray!80}{	<html>}|
		|\colorbox{gray!80}{		<head></head>}|
		|\colorbox{gray!80}{		<body>}|
		|\colorbox{gray!80}{			<p>Hello world</p>}|
		|\colorbox{gray!80}{		<body>}|
		|\colorbox{gray!80}{</html>}|
		|\colorbox{gray!80}{	\}"}|
	\end{minted}
	\caption{A textblock with its actual content highlighted}
\end{listing}

Note how line 16 \code{</html>} now dictates the amount of incidental
whitespaces.

\clearpage

\subparagraph{Indentation leader} Only one type of white space will be
considered as the indentation leader; in case of mixed types, the first
incidental whitespace of the first line (excluding the line with the opening
text block tag) will be considered the indentation leader character. So it is
advised to be consistent with code formatting to avoid ambiguities and
undesired side effects.

\comment{TODO: riferimento a convenzioni e best practices}

\begin{minted}[autogobble=false,showtabs,tab=>,escapeinside=||]{text}
			_out <- "{
		|\colorbox{gray!80}{	<html>}|
		|\colorbox{gray!80}{		<head></head>}|
		|\colorbox{gray!80}{		<body>}|
		|\colorbox{gray!80}{····		<p>Hello world</p>}|
		|\colorbox{gray!80}{		<body>}|
		|\colorbox{gray!80}{	</html>}|
		|\colorbox{gray!80}{	\}"}|
\end{minted}

Note how line 5 accidentally used four spaces instead of a tab; since the tab
was chosen as indentation leader, the spaces are considered non-incidental,
therefore they should be included in the text, hence line 5 now dictates the
amount of incidental whitespaces.

\clearpage

\subsubsection{Numbers}

Giovi supports three different type of numbers: integers, natural and real.

\paragraph{Separators} To improve readability of large numbers, you can use
underscores to group digits; separators cannot appear at the beginning or the
end of the number. For example: \code{100\_000\_000}

\paragraph{Integers} An integer literal may be expressed in multiple
bases.\\*
By default, decimal (base 10) is used; to indicate a different base use the
appropriate prefix:
\begin{itemize}
	\item Decimal (base 10): \code{69420}
	\item Hexadecimal (base 16): \code{0x10F2C}
	\item Octal (base 8): \code{0o207454}
	\item Binary (base 2): \code{0b0001\_0000\_1111\_0010\_1100}
\end{itemize}

\paragraph{Suffixes} By default, a number is parsed as an integer; you can
add a suffix to the number to explicitate the type: \code{n} for naturals,
\code{i} for integers, \code{r} for reals. \code{69n} \code{69i} \code{69r}

\subsubsection{Identifiers}

\comment{TODO: completare}

Any other word that is not a reserved keyword is considered an identifier.

\comment{TODO: restringere i caratteri permessi ad un subset unicode?}

\clearpage

\section{Syntax}
\label{sec:syntax}

This section mainly focuses on the syntax rules of the grammar; for more
details about a specific feature, refer to the dedicated section.

A valid Giovi program starts with zero or more top-level statements; there are
11 valid statements:

\begin{minted}{antlr}
program: statement* ;
statement
	: statement_declaration
	| statement_return
	| statement_invocation
	| statement_redirection
	| statement_import
	| statement_if
	| statement_break
	| statement_skip
	| statement_repeat
	| statement_for
	;

statement_declaration
	: statement_declaration_alias
	| statement_declaration_module
	| statement_declaration_struct
	| statement_declaration_constant
	| statement_declaration_variable
	| statement_declaration_function
	;
\end{minted}

In many rules, a single statement or a block of statements are allowed; the
following rules will be used.

\begin{minted}{antlr}
block: '{' statement* '}';
block_or_statement: block | statement ;
\end{minted}

\paragraph{Identifiers} The \code{identifier} rule describes the allowed
identifiers to be used as names for variables, constants, types etc.

\begin{minted}{antlr}
identifier: [a-zA-Z] [a-zA-Z0-9_]* ;
\end{minted}

It must start with a letter and can contain any letter, number or underscore.

\clearpage

\subsection{Alias declaration}

\begin{minted}{antlr}
statement_declaration_alias: 'alias' expression 'as' identifier ;
\end{minted}

Sometimes a type, a function, a constant or something else may be named in a
too verbose way. You can use the alias statement to basically refer to the same
thing but with a different name.

\begin{minted}{\giovilexer}
alias AVeryVerboseAndLongStructName as SomeStruct
\end{minted}

\subsection{Import statement}

\begin{minted}{antlr}
statement_import
	: 'import' import_entry
	| 'import' '[' import_entry+ ']'
	;
import_entry: import_uri ('.' import_uri)* ('as' alias=identifier)? ;
import_uri
	: identifier
	| literal_normal_text
	| literal_exact_text
	;
\end{minted}

The import statement is used to import one or more modules
(\see{ch:modules}).

\begin{minted}{\giovilexer}
import giovi.util.Random
giovi.util.Random.integer() -> _out
\end{minted}

To avoid fully qualifying the module every time, it can be aliased. It's a
shortcut for importing and then aliasing with two different statements.

\begin{minted}{\giovilexer}
import giovi.util.Random as Random
Random.integer() -> _out
\end{minted}

For more details, see \See{sec:import}.

\subsection{Module declaration}

\begin{minted}{antlr}
statement_declaration_module
	: 'module' identifier body=block
	;
\end{minted}

A module is a "container" for code, and shares the same behavior of a scope: it
can hold code, variables, functions, constants, types definition etcetera; in
fact it can be seen as a "named scope". See \See{ch:modules}.

\nb{Modules are an experimental feature.}

\subsection{Constant declaration}

\begin{minted}{antlr}
statement_declaration_constant
	: 'const' identifier type? '<-' expression
	;
\end{minted}

Constants are read-only values; their value is defined at declaration and
cannot change.\\*
By convention, constants identifiers are written in \code{SCREAMING\_CASE}.\\*
The type can be omitted as it will be inferred from the expression.
\
\begin{minted}{\giovilexer}
const PI Real <- 3.14
const TAU <- PI * 2
\end{minted}

\comment{TODO: decidere se usare shallow o deep immutability}

\subsection{Variable declaration}

\begin{minted}{antlr}
statement_declaration_variable
	: 'var' identifier type? ('default' expression)? '<-' expression
	| identifier type ('default' expression)?
	;
\end{minted}

A variable is identified by a name, and it is of a specific type.\\*
By convention, variable identifiers are written in \code{pascalCase}.\\*
Using the \code{var} keyword allows you to omit the type and let it be inferred
from the assigned expression.

\begin{minted}{\giovilexer}
myName Text <- "Simone"
var myNickname <- "simonvic"
\end{minted}

\subsection{Return statement}

\begin{minted}{antlr}
statement_return
	: ('return' | '=>') expression?
	;
\end{minted}

The return statement can only be used within a function scope.

\subsection{Function declaration}

\begin{minted}{antlr}
statement_declaration_function
	: identifier '(' formal_parameters? ')' ('=>' return_type=type)?
	  body=block
	;

formal_parameters
	: formal_parameter (',' formal_parameter)*
	;

formal_parameter
	: (label=identifier ':')? name=identifier type ('default' expression)?
	;
\end{minted}

Functions are code blocks that can be invoked at a later time.\\*
By convention, function identifiers are written in \code{pascalCase}.

\begin{minted}{\giovilexer}
printHelloWorld() {
	"Hello World" -> _out
}

greetPerson(name Text, surname Text) {
	"Greetings $name, $surname" -> _out
}

getFunLanguage() => Text {
	=> "Giovi"
}
\end{minted}

For more details, see \See{ch:functions}.

\subsection{Invocation statement}

\begin{minted}{antlr}
statement_invocation
	: expression '(' actual_parameters ')'
	;

actual_parameters
	: actual_parameter (',' actual_parameter)*
	;

actual_parameter
	: (label=identifier ':')? expression
	;
\end{minted}

When invoking a function, arguments can be passed using a traditional
positional approach or a nominal approach.\\*
For more details, see \See{ch:functions}.

\nb{Albeit the grammar doesn't enforce one approach or the other, the current
implementation supports either all arguments to be passed as positional or as
nominal.}

\subsection{Redirection statement (assignment)}

\begin{minted}{antlr}
statement_redirection
	: statement_redirection_rtl
	| statement_redirection_ltr
	;

statement_redirection_rtl
	: dest=pattern_redir '<-' src=expression
	| dest=expression ('<-', '+<-', '-<-', '*<-', '/<-') src=expression
	;

statement_redirection_ltr
	: src=expression '->' dest=pattern_redir
	| src=expression ('->', '->+', '->-', '->*', '->/') content=expression
	;

pattern_redir
	: pattern_placeholder
	| pattern_unary
	| pattern_placeholder
	;
\end{minted}

The redirection statement is simply an assignment of a value to
\textit{something}.

A peculiarity of the statement is that it can be written "left to right" or
"right to left". There is no difference between the two directions, except for
readability. Sometimes it's just better to use one instead of the other.

\comment{
	TODO: link a FAQ
}

\begin{minted}{\giovilexer}
"hello world" -> _out     # Printing
counter +<- 1             # Increasing counter by one
acceleration ->+ velocity # Increasing velocity by "acceleration"
\end{minted}

\subsection{If statement}

\begin{minted}{antlr}
statement_if
	: statement_simple_if
	| statement_multi_if
	;
\end{minted}

The \code{if} statement is used to evaluate boolean expressions exclusively.\\*
There is a "traditional" \code{if}, where a single expression is evaluated and
the appropriate execution flow is chosen, and a "multi-if" version.

\subsubsection{Simple If statement}

\begin{minted}{antlr}
statement_simple_if
	: 'if' expression 'do' statement_or_block ('else' statement_or_block)?
	;
\end{minted}

\begin{minted}{\giovilexer}
if someConditionIsTrue() do something() else somethingElse()
if someConditionIsTrue() do {
	something()
} else {
	somethingElse()
}
\end{minted}

\subsubsection{Multi If statement}

\begin{minted}{antlr}
statement_multi_if
	: 'if' subject=expression multi_if_body ('else' statement_or_block)?
	| 'if' subjectless_multi_if_body ('else' statement_or_block)?
	;
\end{minted}

Very often it is needed to check multiple conditions sequentially.

\begin{minted}{\giovilexer}
if name = "Simone" do {
	"hey it's Simone!" -> _out
} else if name = "Mario" do {
	"heyo Mario" -> _out
	"where is Luigi?" -> _out
} else if name = "Luigi" do {
	"heyo green Mario!" -> _out
}
\end{minted}

To improve readability, the \code{if} statement allows you to write the same in
a more concise way.

\begin{minted}{\giovilexer}
if {
	name = "Simone" do "hey it's Simone!" -> _out
	name = "Mario" do {
		"heyo Mario" -> _out
		"where is Luigi?" -> _out
	}
	name = "Luigi" do "heyo green Mario!" -> _out
}
\end{minted}

\nb{Any boolean expression is allowed and exhaustiveness is not enforced! Be
mindful of the branch order; since they're evaluated in the same order they're
written, some branch may become unreachable if it's dominated by a previous
one.}

\subparagraph{Subject} More often than not, the conditions involve the same
"subject". The \code{if} statement allows you to write it just once and then
omit it in the left-hand side of the condition

\begin{minted}{\giovilexer}
if name {
	= "Simone" do "hey it's Simone!" -> _out
	= "Mario" do {
		"heyo Mario" -> _out
		"where is Luigi?" -> _out
	}
	= "Luigi" do "heyo green Mario!" -> _out
}
\end{minted}

Common boolean operators are allowed

\begin{minted}{\giovilexer}
if getSomeNumber() {
	= 69 or = 420 do "Nice!" -> _out
	> 100 and < 999 do "big number" -> _out
}
\end{minted}

A branch can also be "guarded" with another (possibly unrelated) boolean
condition

\begin{minted}{\giovilexer}
var funnyNumbers <- [69, 420] 
if getSomeNumber() {
	# if equal to 69, but only if someCondition() is true
	= 69 when someCondition() do "Nice!" -> _out

	# if is contained in funnyNumbers list,
	# but only when 67 is not contained in funnyNumbers
	in funnyNumbers when 67 not in funnyNumbers
}
\end{minted}

\subsection{Break and Skip statement}

Both statements are used to control the flow of a loop, therefore can only be
used within a loop.

\subsubsection{Break statement}

\begin{minted}{antlr}
statement_break: 'break';
\end{minted}

The break statement will stop the loop and the execution will continue with the
next instruction.

\subsubsection{Skip statement}

\begin{minted}{antlr}
statement_skip: 'skip';
\end{minted}

The skip statement will stop the current loop iteration and skip to the next
iteration.

\subsection{Repeat statement}

\begin{minted}{antlr}
statement_repeat: 'repeat' ('if' expression)? statement_or_block ;
\end{minted}

The repeat statement allows you to repeat a statement or block of code. It
allows an optional condition that gets checked at the beginning of every loop;
if the condition fails, the iteration stops.

\begin{minted}{\giovilexer}
# this code will run indefinitely
repeat {
	"hello world" -> _out
}
\end{minted}

\begin{minted}{\giovilexer}
# the loop continues while someCondition is true
repeat if someCondition() {
	"hello world" -> _out
}
\end{minted}

A "do-while" like statement, can be achieved with a \code{break} statement.

\begin{minted}{\giovilexer}
# the loop continues until someCondition is true
repeat {
	"hello world" -> _out
	if someCondition() do break
}
\end{minted}

\subsection{For statement}

\begin{minted}{antlr}
statement_for
	: 'for' pattern_iterator 'in' iterable 'do' statement_or_block
	;

pattern_iterator: pattern_map_naked | pattern ;
\end{minted}

In its simplest form, the \code{for} statement is sued to iterate over lists,
maps or ranges.

\begin{minted}{\giovilexer}
for i in 0..10 do somethingWith(i)
for number in [1 2 3 4] do somethingWith(number)
for key in ["Simone"=>69 "Luca"=>420 "Paolo"=>10] do somethingWith(key)
\end{minted}

The for statement is like an usual for-each loop, but on steroids; the iterator
is actually a pattern (\see{ch:patterns}): for each iteration, the for
statement will match iterated value with the pattern, and perform the
appropriate deconstruction.

\begin{listing}
	\label{lst:for_loop_patterns_example}
	\caption{Example of various patterns used as iterators in \code{for} loops}
	\begin{minted}{\giovilexer}
for _ in 0..9 do something() # do something 10 times

for [[i]] in some2DMatrix do somethingWith(i)

mapOfNicknames [Text => [Text]] <- [
	"Simone" => ["Simon", "simonvic"]
	"Mario"  => ["Super Mario", "The plumber"]
	"Luigi"  => ["Green Mario"]
]
for [name => [nickname]] in mapOfNicknames do {
	somethingWith(name, nickname)
}

listOfColors [tuple<Int,Int,Int>] <- [
	<0 255 0>
	<55 0 255>
	<255 255 255>
]
for <r, g, b> in listOfColors do {
	"The color is made of $r $g $b" -> _out
}
	\end{minted}
\end{listing}

\nb{Nesting iterator patterns effectively means nesting iterations}

\comment{TODO: aggiungere RFC di iteration counters}

\chapter{Scopes}
\label{ch:scopes}

A scope is a code region within which you can execute statements and perform
declarations (e.g. variables, constants, types, functions) that can later be
referred to, simply, by their name.

A Giovi source code file is itself a scope (see \See{ch:modules}). Scopes can
be "nested" within other scopes. Their lifetime has a clear begin and end; a
scope usually coincides with a \textit{block} (see \See{sec:syntax}), for
example in an \code{if} statement, a \code{for} or \code{repeat} loop or a
function body.

\begin{listing}
	\label{lst:scope_lifetype}
	\caption{Example of nested scopes and their "lifetime"}
	\begin{minted}{\giovilexer}
# begin scope
if true do {
	# begin scope
	for _ in 0..10 do {
		# begin scope
		"hello world" -> _out
		# end scope
	}
	# end scope
} else {
	# begin scope
	"hello world" -> _out
	# end scope
}
#end scope
	\end{minted}
\end{listing}

\clearpage

\section{Declarations lifetime} The lifetime of a declaration is bound to
the lifetime of the scope in which it has been declared.

\begin{minted}{\giovilexer}
if true do {
	username Text <- "simonvic"
	"hello $username" -> _out # "hello simonvic"
}
"hello $username" -> _out # ERROR: username not defined
\end{minted}

A scope "inherits" the declarations of their parent scope; more precisely, a
declaration in a parent scope is visible to the children scopes.

\begin{minted}{\giovilexer}
username Text <- "simonvic"
if true do {
	"hello $username" -> _out # "hello simonvic"
}
\end{minted}

Two declarations of the same type cannot have the same name.

\begin{minted}{\giovilexer}
username Text <- "simonvic"
username Text <- "mario" # ERROR: username already declared
\end{minted}

\section{Shadowing} You can, for example, declare a variable with the same
name of another variable declared in a parent scope; the variable in the parent
scope is now "shadowed", and cannot be referred to by its name.

\begin{minted}{\giovilexer}
username Text <- "simonvic"
if true do {
	username Text <- "mario"
	"hello $username" -> _out # "hello mario"
}
"hello $username" -> _out # "hello simonvic"
\end{minted}

For a more formal definition, see \See{ch:opsem}

\chapter{Defaults}
\label{ch:defaults}

An important design philosophy of the language is that it has no concept (nor
models) a "null" value or uninitialized memory.

Every variable, constant, function parameter, product types fields etc. have a
"default" value, which is the value that they are initialized to at
declaration.

A default value is also referred to as the "zero" value; the two terms will be
used interchangeably.

Every type has a zero value, for example it's the number zero for numeric
types, an empty string for \code{Text}, false for \code{Bool} etc.

\begin{minted}{\giovilexer}
someInteger Int
someInteger -> _out # 0

someText Text
someText -> _out # "" (empty string)

someBool Bool
someBool -> _out # false
\end{minted}

The same applies to lists and maps (empty list and empty map respectively), and
most importantly to product types, where the zero value has all its fields
initialized to their respective zero values.

\begin{listing}
	\label{lst:product_type_zero_value}
	\caption{
		Example of a variable being initialized with the zero value of its
		product type
	}
	\begin{minted}{\giovilexer}
struct Person <
	name Text
	aliases [Text]
	age Int
	isAdult Bool
	favoriteColor tuple<Int, Int, Int>
>
someDefaultPerson Person
someDefaultPerson.name -> _out # ""
someDefaultPerson.aliaes -> _out # []
someDefaultPerson.age -> _out # 0
someDefaultPerson.isAdult -> _out # false
someDefaultPerson.favoriteColor -> _out # <0,0,0>
	\end{minted}
\end{listing}

For a more formal definition, see \See{sec:zero_values}.

\clearpage

\section{The default keyword}

The \code{default} keyword can be used in redirections, for example, to "reset"
a variable to its default value

\begin{minted}{\giovilexer}
age Int
69 -> age
age -> _out # 69
default -> age
age -> _out # 0
\end{minted}

You can get the default value of a type with the dot operator

\begin{minted}{\giovilexer}
Int.default -> _out # 0
Bool.default -> _out # false
Person.default -> _out
# Person.<name:"", aliaes:[], age:0, isAdult:false, favoriteColor:<0,0,0>>
\end{minted}

You can also get the default value of a variable, which will be useful later
with \See{sec:default_override}

\begin{minted}{\giovilexer}
age Int
age.default -> _out # 0
\end{minted}

\section{Overriding default}
\label{sec:default_override}

A peculiarity of the language is that the default value is overrideable.

\nb{This is possible only at declaration.}

\begin{minted}{\giovilexer}
age Int default 69
age.default -> _out # 69
\end{minted}

Product types fields can also have their default value overridden

\begin{listing}
\label{lst:product_type_zero_value_overridden}
	\caption{
		Example of a variable being initialized with the overridden zero value
		of its product type
	}\begin{minted}{\giovilexer}
struct Person <
	name Text default "nobody"
	aliases [Text]
	age Int default 18
	isAdult Bool default true
	favoriteColor tuple<Int, Int, Int> default <255, 255, 255>
>
someDefaultPerson Person
someDefaultPerson.name -> _out # nobody
someDefaultPerson.aliaes -> _out # []
someDefaultPerson.age -> _out # 18
someDefaultPerson.isAdult -> _out # true
someDefaultPerson.favoriteColor -> _out # <255,255,255>
	\end{minted}
\end{listing}

\chapter{Functions}
\label{ch:functions}

\section{Intro}

Functions are code blocks that can be invoked at a later time.

In Giovi, functions are first-class citizens, meaning that they are like any
other value: they have a well defined type (signature), they can be assigned to
variables, returned from other functions or passed to them as arguments, and so
on.

\paragraph{Function signature} A function signature is made of the function
name, its formal parameters (if any) and its return type (if any).

\paragraph{Declaration} A "traditional" function can be declared
by specifying its signature and its body.

\begin{minted}{\giovilexer}
buildGreetingMessage(name Text) => Text {
	if name = "simonvic" do {
		=> "heyo $name, how are you?"
	}
	=> "greetings, $name."
}
\end{minted}

It can later be invoked.

\nb{You have to declare the function first before invoking it. Hoisting is
still an experimental feature.}

\begin{minted}{\giovilexer}
buildGreetingMessage("simonvic") -> _out
\end{minted}

A function can also exist, for example, as a variable, where its type is the
signature. In this case the variable is an "empty function"; before invoking
it, you want to assign an actual function instance (with compatible signature.

\begin{minted}{\giovilexer}
greetingBuilder function(name Text) => Text
greetingBuilder <- buildGreetingMessage
# greetingBuilder is now "referencing" the previously
# defined buildGreetingMessage function
greetingBuilder("simonvic") -> _out
\end{minted}

\section{The dotref operator}

A "traditional" function declaration is semantically similar to a declaration
of a constant of type function, but with a difference: while there cannot be
two equally named functions defined in the same scope, there can be a function
and a variable (or constant) sharing the same name.

However, a function is always shadowed by a variable (or constant) with the
same name. To explicitly refer to the function, the \code{::} prefix is
required.

\begin{minted}{\giovilexer}
var foo <- 69
function foo() { "hello world" -> _out }
foo -> _out # 69
::foo -> _out # function(){}
\end{minted}

The dotref operator can also be used instead of a dot operator

\begin{minted}{\giovilexer}
struct SomeStruct <
	foo Int default 69
	function foo() { "hello world" -> _out }
>

someStruct SomeStruct
someStruct.foo -> _out # 69
someStruct::foo -> _out # function(){}
\end{minted}

\clearpage

\section{Function scope}

The scope of a function behaves like a normal scope, which means that, unlike
many other languages, it doesn't have a "call stack". It can be seen as a
normal scope in which some variables (i.e. formal parameters) are pre-declared
and declarations in parent scopes are visible.

\begin{minted}{\giovilexer}
const DEG_TO_RAD <- 0.01745329
degreeToRadians(deg Int) => Int {
	=> deg * DEG_TO_RAD
}
degreeToRadians(180) -> _out # 3.14
\end{minted}

\nb{This behavior, while it offers a lot of flexibility, it has some caveats:
when a function, that depends on a declaration of a parent scope, is invoked
within a scope where such declaration is not present, the program will error.}

\begin{minted}{\giovilexer}
degreeToRadians(deg Int) => Int {
	=> deg * DEG_TO_RAD
}
if true do {
	const DEG_TO_RAD <- 0.01745329
	degreeToRadians(180) -> _out # 3.14
} else {
	degreeToRadians(180) -> _out # ERROR: DEG_TO_RAD not declared
}
\end{minted}

\subsection{Nested functions}

A function declaration can appear anywhere, even within other function scopes.
Like any other declaration, its lifetime is bound to its enclosing scope,
meaning that after its enclosing scope has ended, the function will not be
visible anymore.

This can be very useful when you have repeating code, but you don't want to
expose a function.

\begin{minted}{\giovilexer}
function makePizza(ingredients [Text]) => Pizza {
	if "pineapple" in ingredients do panic()

	function cook(ingredient Text, seconds Int) {
		# cook the ingredient ...
	}

	cookedIngredients [Text]
	for ingredient in ingredients do {
		if ingredient {
			= "tomato" do cook(ingredient, 60) ->+ cookedIngredients
			= "fries" do cook(ingredient, 30) ->+ cookedIngredients
		} else {
			cook(ingredient, 100) ->+ cookedIngredients
		}
	}
	=> assemblePizza(cookedIngredients)
}
\end{minted}

\clearpage

\section{Parameters and invocation}

A formal parameter is the parameter that is declared in the function signature;
it has a name and a type. A function can have zero or more formal parameters.
They share the variable declaration syntax since they could be seen as such;
when a function is invoked, before executing the first statement of its body, a
new function scope is created and the formal parameters are declared as
variables.

The actual parameters (a.k.a. arguments) of the invocation are passed by value,
which means they're copied, therefore, while the function body can mutate the
parameters, the mutation will not be visible outside the function.

\begin{minted}{\giovilexer}
someFunction(someParameter Int) {
	someParameter -> _out # 69
	someParameter <- 420
	someParameter -> _out # 420
}
someArgument Int <- 69
someFunction(someArgument)
someArgument -> _out # still 69
\end{minted}

\comment{
	TODO: menzionare eccezione per self functions, se decido di "formalizzare
	il bug" :)

	Alternativamente, formalizzare che le strutture (o tutti i prodotti tipi)
	sono mutabili
}

\paragraph{Parameters defaults}

A formal parameter, being like a normal variable, also has a default value. You
can pass a default value by simply using the \code{default} keyword or the
underscore \code{\_}

\begin{minted}{\giovilexer}
someFunction(x Int, y Text) {
	x -> _out # 0
	y -> _out # "" (empty string)
}
someFunction(default, default)
someFunction(_, _)
\end{minted}

\subsection{Optional parameters}

Giovi supports optional parameters: by overriding a formal parameter's default
value, the parameter will be marked as optional. An optional parameter can be
omitted at invocation; its default value will be used instead.

\begin{minted}{\giovilexer}
someFunction(x Int default 69, y Text default "hello world") {
	x -> _out # 69
	y -> _out # hello world
}
someFunction()
\end{minted}

You can mark a parameter as optional without overriding its default

\begin{minted}{\giovilexer}
someFunction(x Int default _) { }
# equivalent to
someFunction(x Int default Int.default) { }
\end{minted}

It's possible to mix optional and required parameters, however only trailing
(last) optional parameters can be omitted when invoking using the positional
approach (\see{sec:arg_approach})

\begin{minted}{\giovilexer}
processPizzaOrder(
  kind Text default "Margherita",
  tableNumber Int,
  withPineapple Bool default false,
) {
	if withPineapple do panic()
	makePizza(kind) -> pizza Pizza
	serve(pizza, tableNumber)
}

processPizzaOrder("Diavola", 69) # "Diavola" to table 69
processPizzaOrder(_, 69) # "Margherita" to table 69
processPizzaOrder(_, _) # "Margherita" to table 0
processPizzaOrder(_, 69, true) # "Margherita" with pineapple to table 69
processPizzaOrder(69) # ERROR: expecting pizza type, but got 69
\end{minted}

\subsection{Positional and nominal parameters}
\label{sec:arg_approach}

\paragraph{Positional parameters} When invoking a function with the positional
approach, you just pass the values as arguments, which will be assigned to the
respective formal parameters in the same order: the first argument is assigned
to the first formal parameter, the second to the second and so on.

\paragraph{Nominal parameters} \label{sec:nominal_params} Giovi supports
nominal arguments: you can specify the name of the argument. When using the
nominal approach, the order of the arguments is not important.

\begin{minted}{\giovilexer}
processPizzaOrder(kind Text, tableNumber Int, withPineapple Bool) { #{...}# }
processPizzaOrder(tableNumber:69, withPineapple:false, kind:"Diavola")
\end{minted}

When using the nominal approach, all optional parameters can be omitted, as
long as the required parameters are specified.

\begin{minted}{\giovilexer}
processPizzaOrder(
  kind Text default "Margherita",
  tableNumber Int,
  withPineapple Bool default false
) { #{...}# }
processPizzaOrder(tableNumber:69) # Margherita to table 69 and no pineapple
\end{minted}

\clearpage

\subsection{Invocation labels}

As we've seen in \See{sec:nominal_params}, parameters names can be specified at
invocation.

When declaring a function, you sometimes have to decide whether to name a
parameter to either favor readability in the function body at the expense of
invocation or viceversa.

% \begin{minted}{\giovilexer}
% processPizzaOrder(
%   kind Text,
%   tableNumber Int,
%   hasPineapple Bool,
% ) {
% 	if hasPineapple do return
% 	makePizza(kind) -> pizza Pizza
% 	serve(pizza, tableNumber)
% }
%
% processPizzaOrder(
%   kind: "Diavola",
%   tableNumber: 69,
%   hasPineapple: false,
% )
% \end{minted}
%
% ... or viceversa
%
% \begin{minted}{\giovilexer}
% processPizzaOrder(
%   ofKind Text
%   atTableNumber Int,
%   withPineapple Bool,
% ) {
% 	if withPineapple do return
% 	makePizza(ofKind) -> pizza Pizza
% 	serve(pizza, atTableNumber)
% }
%
% processPizzaOrder(
%   ofKind: "Diavola",
%   atTableNumber: 69,
%   withPineapple: false,
% )
% \end{minted}

To get the best of both worlds and further improve readability, Giovi allows
you to define an "invocation label". Such label is the name of the parameter
that is exposed to the invocation, while the formal parameter name remains
unchanged in the function body. If an invocation label is not specified, the
formal parameter name will be used.

\begin{minted}{\giovilexer}
processPizzaOrder(
  ofKind: kind Text,
  atTable: tableNumber Int,
  hasPineapple Bool,
) {
	if hasPineapple do panic()
	makePizza(kind) -> pizza Pizza
	serve(pizza, tableNumber)
}

processPizzaOrder(
  ofKind: "Diavola",
  atTable: 69,
  hasPineapple: false
)
\end{minted}

\clearpage

\section{Struct functions}
\label{sec:struct_functions}

Functions can be defined within structs (\see{sec:structs}) and will be
associated to the type.\\*
Beware that these are, by definition, NOT methods; it can be seen as simple
syntax sugar to better organize code using types as "namespaces" for these
functions. We could make an analogy with "static methods" in OO languages.

\begin{minted}{\giovilexer}
Vehicle <
	speed Int

	accelerate(v Vehicle, amount Int) {
		amount ->+ v.speed
	}

	brake(v Vehicle) {
		v.speed -<- 1
	}

	speedToMilesPerHour(v Vehicle) => Int {
		=> v.speed / 1.6093
	}
>

car Vehicle
Vehicle.accelerate(car, 10)
"${car.speed} KMH" -> _out
"${Vehicle.speedToMilesPerHour(car)} MPH" -> _out
\end{minted}

Note that, although functions can be invoked from an instance of a struct,
they're still only related to the type; the syntax is allowed because of self
functions (\see{sec:self_func}).

\subsection{Self functions}
\label{sec:self_func}

More often than not, functions defined within a struct type are meant to
operate on such type (e.g. creating, updating, reading).

The \code{self} keyword can be prefixed to the function declaration; it's a
shortcut for defining a function that has a first parameter of the same type of
the struct in which the function is defined.

\begin{minted}{\giovilexer}
Vehicle <
	speed Int
	self speedToMilesPerHour() => Int {
		=> self.speed / 1.6093
	}
>
\end{minted}

Albeit \code{self} is a reserved keyword and cannot be used as parameter name,
the previous example would be equivalent to the following

\begin{minted}{\giovilexer}
Vehicle <
	speed Int
	speedToMilesPerHour(self Vehicle) => Int {
		=> self.speed / 1.6093
	}
>
\end{minted}

This also allows you to invoke the function using the instance of the struct as
the left-hand side of the dot operator, which will be passed automatically as
the first argument.

\begin{minted}{\giovilexer}
car Vehicle
car.speedToMilesPerHour() -> _out
# equivalent to
Vehicle.speedToMilesPerHour(car) -> _out
\end{minted}

\subsection{Self returning functions}

The \code{self} keyword can also be used as a return type as a shorthand.\\*
This comes very handy for writing "chainable" functions.

\begin{minted}{\giovilexer}
Color <
	r Int
	g Int
	b Int

	rgb(r Int, g Int, b Int) => self {
		=> Color.<r g b>
	}

	self brigthen() => self {
		self.r *<- 0.5
		self.g *<- 0.5
		self.b *<- 0.5
		=> self
	}

	self saturate() => self {
		# ...
		=> self
	}
>

Color.rgb(5, 20, 512)
	.brigthen()
	.saturate()
	.brighten()
\end{minted}

\chapter{Product types}

In this chapter product types will be discussed. A product type is called as
such because it can be seen as a Cartesian product of other types.

There are three product types:
\begin{itemize}
	\item \textbf{Tuple} (\see{sec:tuples})
	\item \textbf{Record} (\see{sec:records})
	\item \textbf{Struct} (\see{sec:structs})
\end{itemize}

All product types follow the same rules of the default (\see{ch:defaults}),
hence, when a variable or any other value of such type is declared, it is
initialized with its zero value; there is one exception (see
\See{sec:recursive_structs}).

\section{Tuples}
\label{sec:tuples}

A Tuple is an ordered collection of zero or more heterogeneous values.

\begin{minted}{\giovilexer}
someTuple tuple<Text, Bool, Real>
someTuple <- <"foo", true, 420.69>
\end{minted}

Tuples being ordered means that two tuple types with the same fields (same
count and equally typed), but in different position, are considered different.

\section{Records}
\label{sec:records}

A Record is an unordered collection of zero or more named (labeled)
heterogeneous values.

\begin{minted}{\giovilexer}
someRecord record<field1 Text, field2 Bool, field3 Real>
someRecord <- <field1:"foo", field2:true, field3:420.69>
\end{minted}

Records being unordered means that two record types with the same fields (same
count, equally named and equally typed) are equal, regardless of the order
their fields appear at.

\section{Structs}
\label{sec:structs}

A Struct is a named unordered collection of zero or more named (labeled)
heterogeneous fields. In other words, a struct type is a type you can define
and later refer to via its name. Structs could be seen as Records
(\see{sec:records}) with a typename associated to it.

Before being able to use a struct type, it must first be declared.

\begin{minted}{\giovilexer}
struct SomeStruct <field1 Text, field2 Bool, field3 Real>
someStruct SomeStruct
\end{minted}

\paragraph{Struct constants} Like for functions (see
\See{sec:struct_functions}), a struct type can also contain constant
declarations, that can later be accessed via the dot operator.

\subsection{Recursive struct types}
\label{sec:recursive_structs}

Recursive struct type definitions are allowed by using the \code{self} keyword
in place of the field type.

\comment{TODO: menzionare RFC di larval types e uso del typename}

\begin{minted}{\giovilexer}
struct Node <
	data Text,
	nextNode self
>
\end{minted}

All types in Giovi have a zero value (see \See{ch:defaults}) and all values
of such type are initialized with it, and the same applies to structs;
however recursive structs are the exception.\\*
In the case of the a recursive struct, the fields of such type will be
initialized with a lazy approach: the initialization of such fields happens
only upon the first access/use.

\begin{minted}{\giovilexer}
struct Node <
	data Text default "default_data",
	nextNode self
>
root Node <- <data:"root_data">
root -> _out       # <"root_data", _>
root.nextNode -> _ # nextNode is now initialized
root -> _out       # <"root_data", <"default_data", _>>
\end{minted}

\subsection{Aggregation and composition}

When modelling a certain domain, it is often needed to define an aggregation or
composition relationship between types. To achieve aggregation, practically
speaking, it's simply possible to define a field of a previously typed.

\begin{minted}{\giovilexer}
struct Fullname <name Text, surname Text>
struct Person <age Int, fullname Fullname>
Person -> _out # Person<age Int, fullname Fullname>
\end{minted}

At times, however, it may be desirable to "flatten/inline" such aggregated
fields; instead of manually copying the fields in the new declaration, it's
possible to specify the typename of the other type you want to aggregate.

\begin{minted}{\giovilexer}
struct Person <age Int> + Fullname
Person -> _out # <age Int, name Text, surname Text>
\end{minted}

The aggregate type can also be specified before the new declaration; by doing
so, the fields will appear before the ones in the new declaration.

\begin{minted}{\giovilexer}
Fullname + struct Person <age Int>
Person -> _out # <name Text, surname Text, age Int>
\end{minted}

Multiple aggregations are possible

\begin{minted}{\giovilexer}
Vehicle, Engine + struct Car<brand Text> + AirCondition, ReclineableSeats
Car -> _out
# Car<
#	Vehicle fields...,
#	Engine fields...,
#	brand Text,
#	AirCondition fields...,
#	ReclineableSeats fields...,
#>
\end{minted}

TODO: clarify multiple aggregation, overrides and defaults

\chapter{Patterns}
\label{ch:patterns}

A pattern is a special language construct that describes a structure of data.

A pattern is used to perform the "deconstruction" of a value into the structure
that the pattern describes.

\begin{minted}{\giovilexer}
getPoint() => tuple<Int, Int> {
	=> <69, 420>
}
<x, y> <- getPoint()
"The point is at $x and $y" -> _out
\end{minted}

Patterns can also be used as iterators in a for loop; at each iteration the
value will be matched to the pattern.

\begin{minted}{\giovilexer}
const CITIES <- [
	"Rome"    => <41.8941047, 12.5151106>
	"NewYork" => <40.7157121, -74.0013879>
	"Tokyo"   => <35.7104635, 139.7300523>
]
for city => <x, y> in CITIES do {
	"$city is located at $x and $y" -> _out
}
\end{minted}

\paragraph{Patterns compatibility} A value can be deconstructed into a pattern
only when it is "equal or narrower" than the pattern; the rules are similar to
types widening rules.

\begin{minted}{\giovilexer}
getPoint() => tuple<Int, Int> { => <69, 420> }
getPoint() -> <x Real, y Real>
"$x and $y are Real" -> _out
\end{minted}

\clearpage

\section{Patterns kinds}

There are five kinds of patterns:
\begin{itemize}
	\item Terminal patterns:
		\begin{itemize}
			\item \textbf{Placeholder pattern}: \code{\_} (underscore)

				\comment{TODO: rinominare. Sinkhole?}
			\item \textbf{Unary pattern}: \code{someIdentifier}
		\end{itemize}
	\item \textbf{Product pattern}:
		\begin{itemize}
			\item \textbf{Positional}: \code{<a, b, c>}
			\item \textbf{Nominal}: \code{<x:a, y:b, z:c>}
		\end{itemize}
	\item Iterator patterns:
		\begin{itemize}
			\item \textbf{List pattern}: \code{[someIdentifier]}
			\item \textbf{Map pattern}: \code{[key => value]}
		\end{itemize}
\end{itemize}

\section{Nesting patterns}

Non-terminal patterns can have other patterns nested.\\*
An extreme example: a value of type \code{[Text=>[tuple<Int,Int,Int>]]} (map
with $\Text$ keys and a list of tuples as value) could be iterated with a
pattern \code{[colorName=>[<r,\_,b>]]} is a map pattern with a unary pattern as
key and a list pattern as value which itself has a positional product pattern
nested, where the second field pattern is a placeholder.

\nb{
	Nesting iterator patterns effectively means nesting iterations.\\*
	For example \code{for [[[i]]] in matrix3D do} would be equivalent of
	writing three nested for loops
}

\subsection{Limitations}

\begin{itemize}
	\item Product patterns can only have terminal patterns (placeholder, unary)
		or other product patterns nested.
	\item In map patterns, the pattern for the key can only be a terminal or a
		product pattern; the pattern for the value can be of any kind (other
		iterator patterns as well).
\end{itemize}

\clearpage

\section{Terminal patterns}

A terminal pattern is called as such because it instructs not to deconstruct
the value further.

\subsection{Placeholder pattern}

The simplest kind of pattern is the placeholder pattern (denoted by an
underscore \code{\_}).\\*
It basically acts as a sinkhole; redirecting to it does nothing. It can be
useful to explicitly discard a function return value, or when nested in other
patterns.

\begin{minted}{\giovilexer}
stack.pop() -> _ # pop a stack but ignore the value

getPoint() -> <_, y>
"I only care about $y" -> _out
\end{minted}

\subsection{Unary pattern}
A unary pattern is a terminal pattern that simply states a single "unit" to
which the value is assigned to as a whole, without deconstruction.\\*
The pattern is denoted by an identifier (the name) and an optional type to
which the deconstructed value is widened to.

\nb{A variable declaration is secretly a unary pattern!}

\clearpage

\section{Product patterns}

A product pattern describes the data structure of a product type, hence it is
meant to be used with product values (tuples, records, structs).

Like mentioned above, the patterns follow the same rules of types
compatibility. In order to be able to deconstruct all product values, you are
given the choice of two deconstruction approaches: positional and nominal.

\subsection{Positional}

With a positional product pattern, the fields of a product value are
deconstructed in the same order they're defined in the pattern.

The positional product pattern can be used to deconstruct all product types.

\begin{minted}{\giovilexer}
<x, y> <- getPoint()
"The point is at $x and $y" -> _out
\end{minted}

\subsection{Nominal}

With the nominal approach you're deconstructing the value based on the fields'
names (the order won't matter), therefore the pattern labels must be the same
of the fields' names.

Obviously, nominal product patterns can only be used with product types with
named fields (i.e. records, structs).

\begin{minted}{\giovilexer}
<y:someY, x:someX> <- getPoint()
"The point is at $someX and $someY" -> _out
\end{minted}

\subsubsection{Implicitly named product pattern}

It is often needed to deconstruct a nominal product value into a pattern whose
variables share the same names.\\*
A shorthand syntax is available for such scenario; however, only unary patterns
are allowed. In this case, the identifier of the unary pattern will be used for
the field lookup of the product value

\begin{minted}{\giovilexer}
<_:y, _:x> <- getPoint()
"The point is at $x and $y" -> _out
\end{minted}

\clearpage

\section{Iterator patterns}

Iterator patters can only be used in a for loop statement (hence the name). For
each iteration of the loop, the iterated value is deconstructed into the
specified pattern.

\subsection{List pattern}

The list pattern is used to iterate over a list; it's basically a unary
pattern, but allows nested patterns.

\begin{minted}{\giovilexer}
# iterate elements of a list
for [i] in [0 1 2 3] do somethingWith(i)

# iterate elements of a nested list
identityMatrix [[Real]] <- [
	[1 0 0]
	[0 1 0]
	[0 0 1]
]
for [[i]] in identityMatrix do somethingWith(i)
\end{minted}

Since it is very common to iterate the elements of a list, the brackets can be
omitted for non-nested patterns

\begin{minted}{\giovilexer}
for i in [0 1 2 3] do somethingWith(i)
\end{minted}

\subsection{Map pattern}

The map pattern is used to iterate over maps; it's a pair of two patterns, one
for the key and one for the value respectively.

\begin{minted}{\giovilexer}
# iterate over a map
for [key => value] in ["Simone"=>69 "Luca"=>420 "Paolo"=>10] do {
	somethingWith(key, value)
}

# iterate over a nested map
graph [Text => [Text => Int]] <- [
	"a" => ["b"=>5 "e"=>3]
	"b" => ["a"=>5 "c"=>1 "d"=>9 "f"=>2]
	"c" => ["b"=>1 "e"=>6]
	"d" => ["b"=>9 "e"=>7]
	"e" => ["a"=>3 "c"=>6 "d"=>7 "f"=>8]
	"f" => ["b"=>2 "e"=>8]
]

for nodeA => [nodeB => edgeCost] in graph do {
	"$nodeA ---$edgeCost---> $nodeB" -> _out
}
\end{minted}

Like for the list pattern, the brackets can be omitted for non-nested patterns

\begin{minted}{\giovilexer}
for key => value in ["Simone"=>69 "Luca"=>420 "Paolo"=>10] do {
	somethingWith(key, value)
}
\end{minted}

\chapter{Modules}
\label{ch:modules}

A module is a "container" for code; it shares the same behavior of scopes
(\see{ch:scopes}): it can hold code, variables, functions, constants, types
definition etcetera; in fact it can be seen as a "named scope".

Each file is implicitly a module, whose name is the filename containing it.

\begin{minted}{\giovilexer}
# my_project/src/Foo.giovi
"hello from Foo module" -> _out
const SOME_CONSTANT Real <- 3.14
someFunction() => Int {
	"hello from Foo.someFunction" -> _out
	=> 69
}
\end{minted}

\section{Inner modules}

A inner sub-module can be declared with the \code{module} keyword.

\begin{minted}{\giovilexer}
# my_project/src/Foo.giovi
"hello from Foo module" -> _out
module Bar {
	"hello from Foo.Bar module" -> _out
	module Baz {
		"hello from Foo.Bar.Baz module" -> _out
	}
}
\end{minted}

\nb{The body of an inner module is not executed until it is imported! See
\See{sec:import}.}

\section{Importing modules}
\label{sec:import}

The \code{import} statement is used to "import" one or more modules.

\begin{minted}{\giovilexer}
module Foo {
	"hello from Foo module" -> _out
}
import Foo # "hello from Foo module"
\end{minted}

\subsection{Import entry}

An import entry can be a qualified identifier (for native and inner modules) or
a file URI for file modules. To later refer to a file module, it must be
aliased first.

\begin{minted}{\giovilexer}
import [
	giovi.util.Math,
	'file://my_project/my_module.giovi' as MyModule
]
giovi.util.Math -> _out
MyModule -> _out
\end{minted}

Import entries URIs are written using the text literals syntax
(\see{sec:text_syntax}) and, as such, it allows for text interpolation, which
is advised to be used carefully.

\begin{minted}{\giovilexer}
nameOfMyModule Text <- "my_other_module"
import "file://my_project/${nameOfMyModule}.giovi"
\end{minted}

It's possible to import to a inner module of a file module with a single import
statement; however the URI can only appear as the first entry

\begin{minted}{\giovilexer}
import 'file://my_project/my_module.giovi'.MySubModule.MySubSubModule
\end{minted}

\subsection{Import behaviour}

The \code{import} keyword may be deceiving as the import statement doesn't
really import a module like in many other languages, but rather it executes its
body and then gives you access to it and its declarations.

\nb{The import of the same module can be performed multiple times, meaning that
its body will be executed for each import.}

\comment{TODO: menzionare funzione util per caching moduli}

\begin{minted}{\giovilexer}
module MyMath {
	"hello from MyMath module" -> _out
	const PI <- 3.14
}
import MyMath # "hello from MyMath module"
MyMath.PI -> _out # 3.14
\end{minted}

\paragraph{Lazy imports} A peculiarity of the import statement is that it can
be performed anywhere like any other statement; the lifetime of the imported
module alias, like all other declarations (e.g. variables, constants,
functions), will be bound to the scope within which the import is executed.\\*
This allows the flexibility of deferring the import to whenever it is actually
needed.

\begin{minted}{\giovilexer}
import giovi.util.Random as Random # immediately imported

if Random.real() > 0.69 do {
	import giovi.util.Math as Math
	Math.PI -> _out # 3.14
}
\end{minted}

\part{Type system}

\chapter{Types}

\section{Conversion}

\resizebox{0.5\linewidth}{!}{
	\begin{tabular}{ll}
		Symbol & Meaning \\\hline
		$=$ & Identity \\\hline
		$w$ & Wideneable \\\hline
		$n$ & Narrowable \\\hline
		$W$ & Wideneable if components are wideneable \\\hline
		$N$ & Narrowable if components are narrowable \\\hline
		$!$ & Explicit cast is required (it may fail) \\\hline
	\end{tabular}
}

\subsection{Primitive types}

\resizebox{0.9\linewidth}{!}{
	\begin{tabular}{|c|c|c|c|c|c|c|c|}
		\hline
		$_{from} \quad ^{to}$ & \Void       & \Type       & \Bool       & \Natural       & \Int       & \Real       & \Text \\\hline
		\Void                 & $=$         & $w$         & $w$         & $w$            & $w$        & $w$         & $w$         \\\hline
		\Type                 &             & $=$         &             &                &            &             &             \\\hline
		\Bool                 &             &             & $=$         &                &            &             &             \\\hline
		\Natural              &             &             &             & $=$            & $w$        & $w$         &             \\\hline
		\Int                  &             &             &             & $n!$           & $=$        & $w$         &             \\\hline
		\Real                 &             &             &             & $n!$           & $n!$       & $=$         &             \\\hline
		\Text                 &             &             &             &                &            &             & $=$         \\\hline
	\end{tabular}
}


\subsection{Non primitive types}

\resizebox{0.9\linewidth}{!}{
	\begin{tabular}{|c|c|c|c|c|c|c|c|}
		\hline
		$_{from} \quad ^{to}$ & \Void       & \List       & \Map        & \Function      & \Tuple       & \Record       & \Struct \\\hline
		\Void                 & $=$         & $w$         & $w$         & $w$            & $w$          & $w$           & $w$           \\\hline
		\List                 &             & $N! W$      &             &                &              &               &               \\\hline
		\Map                  &             &             & $N! W$      &                &              &               &               \\\hline
		\Function             &             &             &             & $N! W$         &              &               &               \\\hline
		\Tuple                &             &             &             &                & $N! W$       & $W$           & $W$          \\\hline
		\Record               &             &             &             &                & $N!$         & $N! W$        & $W$           \\\hline
		\Struct               &             &             &             &                & $N!$         & $N!$          & $N! W!$       \\\hline
	\end{tabular}
}

\chapter{Inference rules}

\section{Notations and Terminology}

\paragraph{Judgements and terms} To describe the type system, its types and
the inference rules, this document uses "judgements" on "terms".\\*
For example, $\tau:T$ is a judgement telling that $\tau$ is a term of type $T$.

\paragraph{Assumptions} Zero or more assumptions (comma separated) can be
used for a judgement; they're specified before the "turnstile" symbol
$\then$.\\*
For example, $\EnvSym \then \tau:T$ can be read as "if $\EnvSym$ is true, then
$\tau$ is a term of type $T$".

\paragraph{Premise} Judgements can be made based on zero or more premises
(semicolon separated judgements).\\*
For example, $\infrule{\EnvSym_1\then\tau_1:T\semiq\EnvSym_2\then\tau_2:T}{
\EnvSym_1,\EnvSym_2\then \code{$\tau_1$ + $\tau_2$}: T}$ means that with a
premise that both judgements $\EnvSym_1\then\tau_1:T$ and
$\EnvSym_2\then\tau_2:T$ are true, then
$\EnvSym_1,\EnvSym_2\then\code{$\tau_1$+$\tau_2$}:T$ is true.

\paragraph{Code} Text written with \code{monospaced} font denotes the text
as written in code verbatim. Some types may also be written with
\texttt{Monospace} font (usually in PascalCase).

If the code evaluates to some value, to denote its type the
"$:\texttt{SomeType}$" notation will be used.\\*
For example: $\code{59 + 10} : \Int$

To denote to what the value evaluates to, the $\eval$ symbol will be used.\\*
For example: $\code{59 + 10} \eval 69$

\paragraph{Functions} Functions are denoted by a list of zero or more input
arguments followed by an output argument.\\*
For example, $(\tau_1,\tau_2)\to\tau_3$ describes a function type that takes
two arguments of type $\tau_1$ and $\tau_2$ respectively, and produces a value
of type $\tau_3$.

\paragraph{"Generic" types} Some types are "generic" (a.k.a. existential
types), meaning that they declare one or more type variable (a.k.a. type
parameters).\\*
The notation $T[\tau:=\sigma]$ describes a generic type $T$ with a type
parameter $\tau$ of some type $\sigma$.\\*
The notation $t : \exists\sigma.T$ means that $t$ is a term of concrete
type $T$ with the type parameter $\tau$ equal to $\sigma$.

\comment{TODO: tipi esistenziali? Rivedere notazione}

\paragraph{List} To denote a list type, the existential type notation will
be used $\List[\nu=\Type]$.\\*
To denote a value, for example of type $\exists\Int.\List$, this notation will
be used: $\glist{1,2,3}$

\paragraph{Map} To denote a map type, the existential type notation will
be used $\Map[\kappa:=\Type,\nu:=\Type]$.\\*
To denote a value, for example of type $\exists\Text.\exists\Int.\Map$, this
notation will be used: $\map{a\mapsto1,b\mapsto2,c\mapsto3}$

\paragraph{Identifiers set} A set $\Id=\set{i\mid i:=\text{identifier}}$ is
used to encompass all possible identifiers allowed by the grammar.

\comment{TODO: formalizzare}

\paragraph{Product types}

\subparagraph{Tuple} A tuple type is denoted with the cartesian product
operator.\\*
For example, $\tau_1 \times \tau_2 \times \tau_3$ describes a 3-tuple type
with three fields of type $\tau_1, \tau_2, \tau_3$.

\subparagraph{Record} The notation $\tuple{x_1:\tau_1, x_2:\tau_2}$ describes a
record type (a tuple whose fields are named).

\subparagraph{Struct} The notation $(\xi, \tuple{x_1:\tau_1, x_2:\tau_2})$
describes a struct type (a pair of an identifier $\xi$ and a record type).

\comment{TODO: rivedere notazione struct}

\begin{definition}[Types widening/narrowing]
	A type $N$ is said to be "narrower" than a type $W$ iff a value described
	by $N$ can also be described by $W$ without any loss of information.\\*
	Conversely, a type $W$ is said to be "wider" than a type $N$ iff it can
	describe any value described by $N$ without loss of information.\\*
	A "narrower than" relationship is denoted using the $\narreq$ operator.
	For example: $N \narreq W$.
\end{definition}

The $\sqcap$ operator produces the least wide type for the operand types
\[
	\infrule[aligned]{
		\EnvSym & \then \tau_1:\Type \semiq \tau_2:\Type \\*
		\EnvSym & \then \exists_{\omega:\Type}{\tau_1\narreq\omega \land \tau_2\narreq\omega} \\*
		\EnvSym & \then \nexists_{\omega'\Type}{\tau_1\narreq\omega' \land \tau_2\narreq\omega' \land \omega'\narreq\omega}
	}{
		\EnvSym\then \tau_1 \sqcap \tau_2 = \omega
	}
\]

\clearpage

\subsection{Type environments}

The definition of an expression type is done within a "context" (a.k.a. type
environment) which describes the type of any identifier or any other
expressions appearing in the described expression.

The general type environment is denoted by $\EnvSym$, a function that maps an
identifier to its type.

To avoid ambiguities and to more clearly describe expressions, $\EnvSym$ is
partioned in the $\EnvVar,\EnvConst,\EnvFunc,\EnvStucts$ environments:

\begin{itemize}
	\item $\EnvVar$: the Variable environment, a function that maps a variable
		identifier $i$ to its type $\tau$
		\[
			\begin{split}
				\EnvVar : \Id &\to \Type\\*
				\EnvVar(i)    &\mapsto \tau
			\end{split}
		\]
	\item $\EnvConst$: the Constants environment, a function that maps a
		constant identifier $i$ to its type $\tau$
		\[
			\begin{split}
				\EnvConst : \Id &\to \Type\\*
				\EnvConst(i)    &\mapsto \tau
			\end{split}
		\]
	\item $\EnvFunc$: the Functions environment, a function that maps a
		function identifier $i$ to its definition signature $(i, \phi)$
		\[
			\begin{split}
				\EnvFunc : \Id &\to (\Id, \Function)\\*
				\EnvFunc(i)    &\mapsto (i, \phi)
			\end{split}
		\]
		With a type $\tau$, it maps to the function defined within the type
		$\tau$
		\[
			\begin{split}
				\EnvFunc : (\Type,\Id) &\to (\Id, \Function)\\*
				\EnvFunc(\tau, i)      &\mapsto (i, \phi)
			\end{split}
		\]
	\item $\EnvStucts$: the Structs environment
		\[
			\begin{split}
				\EnvStucts : \Id &\to (\Id, \Record)\\*
				\EnvStucts(i)    &\mapsto (i, \rho)
			\end{split}
		\]
		A function that maps a typename $i$ (struct definition identifier) to
		its type definition (identifier $i$ and record $\rho$)
	% \item $\EnvModules$: the Modules environment
	% 	\[
	% 		\begin{split}
	% 			\EnvModules : \Id & \to \Module\\*
	% 			\EnvModules(i)    & \mapsto m
	% 		\end{split}
	% 	\]
	% 	A function that maps an identifier $i$ to the module
\end{itemize}

\comment{TODO: formalizzare structs, funzioni e moduli}

The $\EnvSym$ function can now be defined as follows
\[
	\begin{split}
		\EnvSym : \Id & \to \Type\\*
		\EnvSym(i) & \mapsto \begin{cases}
			\EnvStucts(i) &\text{ if } \EnvVar(i)=\varnothing \land \EnvConst(i)=\varnothing \land \EnvFunc(i)=\varnothing \\*
			\EnvFunc(i) &\text{ if } \EnvVar(i)=\varnothing \land \EnvConst(i)=\varnothing \\*
			\EnvConst(i) &\text{ if } \EnvVar(i)=\varnothing \\*
			\EnvVar(i) & \\*
			% \EnvModules(i) & \text{ if } \\*
		\end{cases}
	\end{split}
\]
which formally describes the lookup order and priority of a symbol $i$.

\paragraph{Environment extension} Given how scopes and symbols lookup work, it
will be useful to define "new" type environments extending an existing
environment. The notation $E\env{t:\tau}$ will be used as a shorthand to denote
$E \cup \env{t:\tau}$, which is a new environment that inherits from $E$ and
extends it with $\env{t:\tau}$.

Using the $\EnvVar$ variables environment as example
\[
	\EnvVar\env{t:\tau}(u) \mapsto \begin{cases}
		T &\text{ if } u = t \\*
		\EnvVar(u) & \\*
	\end{cases}
\]

\clearpage

\section{Types}

\nb{
	The syntax for expressing a type or an expression may be ambiguous without
	context. For this reason the \code{type.} prefix syntax is used, but note
	that in some contexts (e.g. variable declaration) the \code{type.} is not
	needed.
}

\comment{TODO: menzionare uso di type system strutturale e nominale}

\comment{TODO: usare 0-based indexes}

\comment{TODO: usare notazione tipi polimorfici invece che solo Type}

\subsection{Primitive types}

\[
	\begin{split}
		&\EnvSym\then \Void:\Type \\
		&\EnvSym\then \Bool:\Type \\
		&\EnvSym\then \Natural:\Type \\
		&\EnvSym\then \Int:\Type \\
		&\EnvSym\then \Real:\Type \\
		&\EnvSym\then \Text:\Type \\
	\end{split}
\]

\subsection{Collection types}

\[
	\begin{split}
		&\EnvSym\then \List[\nu=\Type]:\Type \\
		&\EnvSym\then \Map[\kappa:=\Type,\nu:=\Type]:\Type \\
	\end{split}
\]

\[
	\infrule{
		\EnvSym\then \tau:\Type
	}{
		\EnvSym\then \code{type.[$\tau$]} \eval \exists\tau.\List
	}
\]

\[
	\infrule{
		\EnvSym\then \tau:\Type \semiq \sigma:\Type
	}{
		\EnvSym\then \code{type.[$\tau$ => $\sigma$]} \eval \exists\tau.\exists\sigma.\Map
	}
\]

\nb{
	When the type parameters are not needed in the rule, they will be omitted.
	For example: $\tau=\List$, $\expr:\Map$
}

\subsection{Product types}

\paragraph{Tuple}

\[
	\EnvSym\then \Tuple : \Type
\]

\nb{
	When the whole definition of a tuple type is not needed, the type will
	simply be denoted by \Tuple.
}

\[
	\infrule{
		\EnvSym\then n \geq 0 \semiq \forall_{i\in[1,n]}{\tau_i:\Type}
	}{
		\EnvSym\then \code{tuple <$\tau_1$, $\dotsc$, $\tau_n$>} \eval \tau_1 \times \dotsb \times \tau_n
	}
\]

\paragraph{Record}

\[
	\EnvSym\then \Record : \Type
\]

\nb{
	When the whole definition of a record type is not needed, the type will
	simply be denoted by \Record.
}

\[
	\infrule{
		\EnvSym\then n \geq 0 \semiq \forall_{i\in[1,n]}{\tau_i:\Type, x_i \in \Id}
	}{
		\EnvSym\then \code{record <$x_1$ $\tau_1$, $\dotsc$, $x_n$ $\tau_n$>} \eval \tuple{x_1:\tau_1, \dotsc, x_n:\tau_n}
	}
\]

\paragraph{Struct}

\[
	\EnvSym\then \Struct : \Type
\]

\nb{
	When the whole definition of a struct type is not needed, the type will
	simply be denoted by \Struct.
}

\[
	\infrule{
		\EnvSym\then \xi \in \Id \semiq n \geq 0 \semiq \forall_{i\in[1,n]}{\tau_i:\Type, x_i \in \Id}
	}{
		\EnvSym,\code{struct $\xi$ <$x_1$ $\tau_1$, $\dotsc$, $x_n$ $\tau_n$>} \then \code{$\xi$} \eval (\xi,\tuple{x_1:\tau_1,\dotsc,x_n:\tau_n})
	}
\]

\comment{TODO: che notazione usare per \Struct???}

\subsection{Functions}

\[
	\EnvSym\then \Function : \Type
\]

\nb{
	When the whole definition of a function type is not needed, the type will
	simply be denoted by \Function.
}

\[
	\infrule{
		\EnvSym\then n \geq 0 \semiq \forall_{i\in[1,n]}{\tau_i:\Type,x_i\in\Id} \semiq \tau_{n+1}\Type
	}{
		\EnvSym\then \code{::($x_1$ $\tau_1$, $\dotsc$, $x_n$ $\tau_n$) => $\tau_{n+1}$} \eval (\tau_1,\dotsc,\tau_n)\to\tau_{n+1}
	}
\]

\[
	\infrule{
		\EnvSym\then n \geq 0 \semiq \forall_{i\in[1,n]}{\tau_i:\Type,x_i\in\Id}
	}{
		\EnvSym\then \code{::($x_1$ $\tau_1$, $\dotsc$, $x_n$ $\tau_n$)} \eval (\tau_1,\dotsc,\tau_n)\to\Void
	}
\]

\[
	\infrule{
		\EnvSym\then \xi \in \Id \semiq n \geq 0 \semiq \forall_{i\in[1,n]}{\tau_i:\Type,x_i\in\Id} \semiq \tau_{n+1}\Type
	}{
		\EnvSym,\code{$\xi$($x_1$ $\tau_1$, $\dotsc$, $x_n$ $\tau_n$) => $\tau_{n+1}$ \{\}} \then \code{$\xi$} : (\tau_1,\dotsc,\tau_n)\to\tau_{n+1}
	}
\]

\[
	\infrule{
		\EnvSym\then \xi \in \Id \semiq n \geq 0 \semiq \forall_{i\in[1,n]}{\tau_i:\Type,x_i\in\Id}
	}{
		\EnvSym,\code{$\xi$($x_1$ $\tau_1$, $\dotsc$, $x_n$ $\tau_n$) \{\}} \then \code{$\xi$} : (\tau_1,\dotsc,\tau_n)\to\Void
	}
\]

\subsection{Range}

\comment{TODO: tipo concreto non implementato ancora}

\subsection{Compatibility rules}

\comment{TODO: menzionare covarianze (tipi composti e prodotto) e
controvarianza (parametri funzioni?)}

\paragraph{Reflexivity}
\[
	\infrule{
		\EnvSym\then \tau:\Type
	}{
		\EnvSym\then \tau \narreq \tau
	}
\]

\paragraph{Transitivity}
\[
	\infrule{
		\EnvSym\then \tau_1 \narreq \tau_2 \semiq \tau_2 \narreq \tau_3
	}{
		\EnvSym\then \tau_1 \narreq \tau_3
	}
\]

\subsubsection{Primitive types}

\[
	\infrule{
		\EnvSym\then \tau : \Type
	}{
		\EnvSym\then \Void \narreq \tau
	}
\]

\[
	\EnvSym\then \Natural \narreq \Int
	\qquad \EnvSym\then \Int \narreq \Real
\]

\subsubsection{Collection types}

\[
	\infrule{
		\EnvSym\then \tau_1 = \List[\nu_1:=\Type]\\*
		\EnvSym\then \tau_2 = \List[\nu_2:=\Type]\\*
		\EnvSym\then \nu_1 \narreq \nu_2
	}{
		\EnvSym\then \tau_1 \narreq \tau_2
	}
\]

\[
	\infrule{
		\EnvSym\then \tau_1 = \Map[\kappa_1:=\Type, \nu_1:=\Type]\\*
		\EnvSym\then \tau_2 = \Map[\kappa_2:=\Type, \nu_2:=\Type]\\*
		\EnvSym\then \kappa_1 \narreq \kappa_2 \land \nu_1 \narreq \nu_2
	}{
		\EnvSym\then \tau_1 \narreq \tau_2
	}
\]

\clearpage

\subsubsection{Product types}

Since the widening and narrowing is crucial for type inference, Tuples, Records
and Structs require some additional attention as they have some differences.

\vspace{\baselineskip}

In order to quickly understand why the widening between product types behaves
the way it does, ask yourself: \textit{what information does this data
carry?}\\*
For example, \code{<1,2,3,"foo">} is a tuple with three integer fields and a
text field; no additional semantic information is given, and those values could
mean anything.\\*
Instead, \code{<x:1, y:2, z:3, name:"foo">} is a record with three integer
fields and a text field, with the additional information about their name;
their semantic meaning is still uncertain.\\*
Finally, \code{Rgb.<x:1, y:2, z:3, name:"foo">} is a struct describing a color
which has an integer field for each red, green and blue channel, and a field
representing the common name of the color.\\*
Instead, \code{Point.<x:1, y:2, z:3, name:"foo">} is a struct describing
a point in a coordinate system which has a three integer field for each
dimension and a label.

\vspace{\baselineskip}

While all previous examples evaluates to an "equal" value (the data they
carry), their semantic meaning is potentially very different. This means that,
while the tuple or the record could be widened to either \code{Rgb} or
\code{Point} (where the "added information" is their semantic meaning), a
\code{Point} cannot be losslessly converted to \code{Rgb} or other narrower
types, because we would "lose" their semantic meaning.\\*
If you're ok with losing such information, conversion between "incompatible"
product types can still be achieved via explicit type coercion (casting).

\clearpage

\paragraph{Tuple}

\[
	\infrule{
		\EnvSym\then n \geq 0 \semiq \forall_{i\in[1,n]}{\tau_i:\Type} \semiq \theta_1 = \bigtimes_{i \in [1,n]}{\tau_i}\\*
		\EnvSym\then m \geq n \semiq \forall_{i\in[1,m]}{\omega_i:\Type} \semiq \theta_2 = \bigtimes_{j \in [1,m]}{\omega_j}\\*
		\EnvSym\then \forall_{i \in [1,n]}{\tau_i \narreq \omega_i}
	}{
		\EnvSym\then \theta_1 \narreq \theta_2
	}
\]
A tuple type $\theta_1$ is narrower than the tuple type $\theta_2$ if the first
$n$ fields of $\theta_1$ have a type $\tau$ narrower than the first $n$ fields
types $\omega$ of $\theta_2$.

\vspace{\baselineskip}

\[
	\infrule{
		\EnvSym\then n \geq 0 \semiq \forall_{i\in[1,n]}{\tau_i:\Type} \semiq \theta = \bigtimes_{i \in [1,n]}{\tau_i}\\*
		\EnvSym\then m \geq n \semiq \forall_{j\in[1,m]}{x_j \in \Id, \omega_j:\Type} \semiq \rho = \tuple{x_1:\omega_1,\dotsc,x_m:\omega_m} \\*
		\EnvSym\then \forall_{i \in [1,n]}{\tau_i \narreq \omega_i}
	}{
		\EnvSym\then \theta \narreq \rho
	}
\]
A tuple type $\theta$ is narrower than the record type $\rho$ if the first $n$
fields of $\theta$ have a type $\tau$ narrower than the first $n$ fields type
$\omega$ of $\rho$ (fields names are ignored).

\vspace{\baselineskip}

\[
	\infrule{
		\EnvSym\then n \geq 0 \semiq \forall_{i\in[1,n]}{\tau_i:\Type} \semiq \theta = \bigtimes_{i \in [1,n]}{\tau_i}\\*
		\EnvSym\then \xi \in \Id \semi m \geq n \semi \forall_{j\in[1,m]}{x_j \in \Id, \omega_j:\Type} \semi \sigma = (\xi, \tuple{x_1:\omega_1,\dotsc,x_m:\omega_m}) \\*
		\EnvSym\then \forall_{i \in [1,n]}{\tau_i \narreq \omega_i}
	}{
		\EnvSym\then \theta \narreq \sigma
	}
\]
A tuple type $\theta$ is narrower than the struct type $\sigma$ if the first
$n$ fields of $\theta$ have a type $\tau$ narrower than the first $n$ fields
type $\omega$ of $\sigma$ (struct name and fields names are ignored).

\vspace{\baselineskip}

\clearpage

\paragraph{Record}

\[
	\infrule{
		\EnvSym\then n \geq 0 \semiq \forall_{i\in[1,n]}{x_i \in \mathcal{N} \subseteq \mathcal{M}, \tau_i:\Type} \semiq \rho_1 = \tuple{x_1:\tau_1,\dotsc,x_n:\tau_n} \\*
		\EnvSym\then m \geq n \semiq \forall_{j\in[1,m]}{y_j \in \mathcal{M} \subseteq \Id, \omega_j:\Type} \semiq \rho_2 = \tuple{y_1:\omega_1,\dotsc,y_m:\omega_m} \\*
		\EnvSym\then \forall_{x':\tau'\in\mathcal{N}}{\exists_{y':\omega'\in\mathcal{M}}{x'=y' \land \tau' \narreq \omega'}}
	}{
		\EnvSym\then \rho_1 \narreq \rho_2
	}
\]
A record type $\rho_1$ is narrower than a record type $\rho_2$ if, for all
fields $x'$ of type $\tau'$ in $\rho_1$, there is a field with the same name
$y'$ of type $\omega'$ in $\rho_2$, where $\tau'$ is narrower than $\omega'$.

\vspace{\baselineskip}

\[
	\infrule{
		\EnvSym\then n \geq 0 \semiq \forall_{i\in[1,n]}{x \in \mathcal{N} \subseteq \mathcal{M}, \tau_i:\Type} \semiq \rho = \tuple{x_1:\tau_1,\dotsc,x_n:\tau_n} \\*
		\EnvSym\then \xi \in \Id \semi m \geq n \semi \forall_{j\in[1,m]}{y_j \in \Id, \omega_j:\Type} \semiq \sigma = (\xi, \tuple{y_1:\omega_1,\dotsc,y_m:\omega_m}) \\*
		\EnvSym\then \forall_{x':\tau'\in\mathcal{N}}{\exists_{y':\omega'\in\mathcal{M}}{x'=y' \land \tau' \narreq \omega'}}
	}{
		\EnvSym\then \rho \narreq \sigma
	}
\]
A record type $\rho$ is narrower than a struct type $\sigma$ if, for all fields
$x'$ of type $\tau'$ in $\rho$, there is a field with the same name $y'$ of
type $\omega'$ in $\sigma$, where $\tau'$ is narrower than $\omega'$ (struct
name is ignored).

\paragraph{Struct}

\[
	\infrule{
		\EnvSym\then \xi_1 \in \Id \semiq n \geq 0 \semiq \forall_{i\in[1,n]}{x_i \in \Id, \tau_i:\Type} \\*
		\EnvSym\then \sigma_1 = (\xi_1, \tuple{x_1:\tau_1,\dotsc,x_n:\tau_n}) \\*
		\EnvSym\then \xi_2 \in \Id \semiq m \geq n \semiq \forall_{j\in[1,m]}{y_j \in \Id, \omega_j:\Type} \\*
		\EnvSym\then \sigma_2 = (\xi_2, \tuple{y_1:\omega_1,\dotsc,y_m:\omega_m}) \\*
		\EnvSym\then \xi_1 = \xi_2
	}{
		\EnvSym\then \sigma_1 \narreq \sigma_2
	}
\]
A struct type $\sigma_1$ is narrower than a struct type $\sigma_2$ if and only
if they have the same name (fields are ignored since they're equal by
construction).

\comment{TODO: questa regola è ridondante?}
\vspace{\baselineskip}

\clearpage

\subsubsection{Function}

\[
	\infrule{
		\EnvSym\then n \geq 0 \semiq \forall_{i\in[1,n]}{x_i\in\mathcal{N}\subseteq\mathcal{M},\tau_i:\Type} \semiq \tau_{n+1}:\Type  \\*
		\EnvSym\then \phi_1 = (x_1:\tau_1,\dotsc,x_n:\tau_n)\to\tau_{n+1} \\*
		\EnvSym\then m \geq n \semiq \forall_{j\in[1,m]}{y_j\in\mathcal{M}\subseteq\Id,\omega_i:\Type} \semiq \omega_{m+1}:\Type \\*
		\EnvSym\then \phi_2 = (y_1:\omega_1,\dotsc,y_m:\omega_m)\to\omega_{m+1} \\*
		\EnvSym\then \tau_{n+1} \narreq \omega_{m+1} \\*
		\EnvSym\then \forall_{x':\tau'}{\exists_{y':\omega'}{x'=y' \land \tau' \narreq \omega'}}
	}{
		\EnvSym\then \phi_1 \narreq \phi_2
	}
\]
A function type $\phi_1$ is narrower than a function type $\phi_2$ if the
return type $\tau_{n+1}$ is narrower than the other return type $\omega_{m+1}$,
and if, for all parameters $x'$ of type $\tau'$ in $\phi_1$, there is a
parameter with the same name $y'$ of type $\omega'$ in $\phi_2$, where $\tau'$
is narrower than $\omega'$.

\comment{
	TODO: è gia un casino cosi :D, ma includere anche opzionalità dei parametri\\
	Eventualmente estrarre una regola per il singolo parametro
}


\commentinline{
	TODO: come menzionato nel README, il widening tra funzioni potrebbe essere
	errato se si considerano i domini delle funzioni rispetto al loro corpo.

	Per funzioni senza parametri è ok: \code{::() => Int} può essere allargata
	a \code{::() => Real}.

	Ma se una funziona richiede dei parametri \code{::(i Int) => Int},
	significa che il suo corpo sa come lavorare con \code{i Int}; nel momento
	in cui si allarga a \code{::(i Real) => Real}, il corpo potrebbe operare su
	un tipo inaspettato
}

\clearpage

\section{Expressions}

\subsection{Literals}

\subsubsection{Bool}

\[
	\EnvSym\then \code{true}:\Bool \qquad \EnvSym \then \code{false}:\Bool
\]

\subsubsection{Int}

\[
	\infrule{
		\EnvSym\then x \in \mathbb{Z} \semiq x \in [-2^{31}, 2^{31} - 1]
	}{
		\EnvSym\then \code{$x$} : \Int \\*
		\land \EnvSym\then \code{$x$i} : \Int
	}
\]

% \subparagraph{Examples} \verb|0, 0i, 69, -420, 32i, -99i|

\subsubsection{Natural}

\[
	\infrule{
		\EnvSym\then x \in \mathbb{N} \semiq x \in [0, 2^{31}]
	}{
		\EnvSym\then \code{$x$n} : \Natural
	}
\]

% \subparagraph{Examples} \verb|0n, 69n, 420n|

\subsubsection{Real}

\[
	\infrule{
		\EnvSym\then x \in \mathbb{R}
	}{
		\EnvSym\then \code{$x$} : \Real \\*
		\land \EnvSym\then \code{$x$r} : \Real
	}
\]

% \subparagraph{Examples} \verb|0r, 0.0, 6.9, -4.20|

\subsubsection{Text}

\[
	\infrule{
		\EnvSym\then x \in \text{Unicode} \setminus \set{"\backslash n", "\backslash r"} \semiq n \geq 0
	}{
		\EnvSym\then \code{"$x_1 \dotsc x_n$"} : \Text \\*
		\land \EnvSym\then \code{'$x_1 \dotsc x_n$'} : \Text
	}
\]

\[
	\infrule{
		\EnvSym\then x \in \text{Unicode} \semiq n \geq 0
	}{
		\EnvSym\then \code{"\{$x_1 \dotsc x_n$\}"} : \Text \\*
		\land \EnvSym\then \code{'\{$x_1 \dotsc x_n$\}'} : \Text
	}
\]

% \subparagraph{Examples} \verb|"hello", 'world', "{hello}", '{world}'|

\comment{TODO: specificare il set $Unicode$ con i caratteri permessi da Java}

\clearpage

\subsection{Collection types}

\subsubsection{List}

\[
	\EnvSym\then \code{[]}:\exists\Void.\List
\]

\[
	\infrule{
		\EnvSym\then n \geq 1 \semiq \tau:\Type \semiq \forall_{i\in[1,n]}{\expr_i\in\Expr,\expr_i:\tau}
	}{
		\EnvSym\then \code{[$\expr_1$, $\dotsc$, $\expr_n$]}:\exists\tau.\List
	}
\]

\[
	\infrule{
		\EnvSym\then n \geq 1 \semiq \forall_{i\in[1,n]}{\tau_i:\Type,\expr_i\in\Expr,\expr_i:\tau_i} \semiq \omega = \bigsqcap_{i=1}^{n}{\tau_i}
	}{
		\EnvSym\then \code{[$\expr_1$, $\dotsc$, $\expr_n$]}:\exists\omega.\List
	}
\]

List of $\omega$ (widest among $\tau$)

\subsubsection{Map}

\[
	\EnvSym\then \code{[=>]}:\exists\Void.\exists\Void.\Map
\]

\[
	\infrule{
		\EnvSym\then \kappa:\Type \semiq \forall_{i\in[1,n]}{k_i\in\Expr,k_i:\kappa} \semiq n \geq 1 \\*
		\EnvSym\then \nu:\Type \semiq \forall_{i\in[1,n]}{v_i\in\Expr,v_i:\nu} \hfill
	}{
		\EnvSym\then \code{[$k_1$=>$v_1$, $\dotsc$, $k_n$=>$v_n$]}:\exists\kappa.\exists\nu.\Map
	}
\]

\[
	\infrule{
		\EnvSym\then \forall_{i \in [1,n]}{\kappa_i:\Type, k_i\in\Expr, k_i:\kappa_i} \semiq n \geq 1 \\*
		\EnvSym\then \forall_{i \in [1,n]}{\nu_i:\Type, v_i\in\Expr, v_i:\nu_i} \hfill\\*
		\EnvSym\then \omega_\kappa = \bigsqcap_{i=1}^{n}{\kappa_i} \semiq \omega_\nu = \bigsqcap_{i=1}^{n}{\nu_i}
	}{
		\EnvSym\then \code{[$k_1$=>$v_1$, $\dotsc$, $k_n$=>$v_n$]}:\exists\omega_\kappa.\exists\omega_\nu.\Map
	}
\]

Map of $\omega_\kappa$ keys (widest  among $\kappa$) and $\omega_\nu$ (widest
among $\nu$).

\clearpage

\subsection{Product types}

\subsubsection{Tuple}

\[
	\infrule{
		\EnvSym\then n \geq 0 \semiq \forall_{i\in[1,n]}{\tau_i:\Type, \expr_i\in\Expr, \expr_i:\tau_i}
	}{
		\EnvSym\then \code{<$\expr_1$, $\dotsc$, $\expr_n$>}:\tau_1 \times \dotsc \times \tau_n
	}
\]

\subsubsection{Record}

\[
	\infrule{
		\EnvSym\then n \geq 1 \semiq \forall_{i\in[1,n]}{x_i\in\Id, \tau_i:\Type, \expr_i\in\Expr, \expr_i:\tau_i}
	}{
		\EnvSym\then \code{<$x_1$:$\expr_1$, $\dotsc$, $x_n$:$\expr_n$>}:\tuple{x_1:\tau_1,\dotsc,x_2:\tau_n}
	}
\]

\subsubsection{Struct}

\[
	\infrule{
		\EnvSym\then \xi \in \Id \semiq n \geq 0 \semiq \forall_{i\in[1,n]}{x_i \in \Id, \tau_i:\Type} \semiq \sigma = (\xi, \tuple{x_1:\tau_1,\dotsc,x_n:\tau_n}) \\*
	}{
		\EnvSym\then \code{$\xi$.<>} : \sigma \\*
	}
\]
This expression initializes a struct $\xi$ of type $\sigma$.

\[
	\infrule{
		\EnvSym\then \xi \in \Id \semiq n \geq 0 \semiq \forall_{i\in[1,n]}{x_i \in \Id, \tau_i:\Type} \semiq \sigma = (\xi, \tuple{x_1:\tau_1,\dotsc,x_n:\tau_n}) \\*
		\EnvSym\then \forall_{i\in[1,n]}{\eta_i:\Type,\eta_i\narreq\tau_i,\expr_i:\eta_i} \\*
	}{
		\EnvSym\then \code{$\xi$.<$\expr_1$, $\dotsc$, $\expr_n$>} : \sigma \\*
	}
\]
This expression initializes a struct $\xi$ of type $\sigma$ by using a
tuple-like notation, where each field gets initialized with the expression
$\expr_i$ of type $\eta_i$ (narrower than $\tau_i$).

\[
	\infrule{
		\EnvSym\then \xi \in \Id \semi n \geq 0 \semi \forall_{i\in[1,n]}{x_i\in\mathcal{N}\subseteq\Id, \tau_i:\Type} \semi \sigma = (\xi, \tuple{x_1:\tau_1,\dotsc,x_n:\tau_n}) \\*
		\EnvSym\then m \leq n \semi \forall_{j\in[1,m]}{y_j\in\mathcal{M}\subseteq\mathcal{N}, \eta_j:\Type, \expr_j:\eta_j} \\*
		\EnvSym\then \forall_{j\in[1,m]}{\exists_{i\in[1,n]}{y_j=x_i \land \eta_j\narreq\tau_i}}
		% \EnvSym\then \forall_{\mathcal{M}\ni y':\eta'}{\exists_{\mathcal{N}\ni x':\tau'}{y'=x' \land \eta'\narr\tau'}}
	}{
		\EnvSym\then \code{$\xi$.<$y_1$:$\expr_1$, $\dotsc$, $y_n$:$\expr_n$>} : \sigma \\*
	}
\]
This expression initializes a struct $\xi$ of type $\sigma$ by using a
record-like notation, where each specified field $y_j$ gets initialized with the
expression $\expr_j$ of type $\eta_j$ (narrower than $\tau_i$).

\clearpage

\subsection{Parenthesized expression}

\[
	\infrule{
		\EnvSym\then \tau:\Type \semiq \expr\in\Expr \semiq \expr:\tau
	}{
		\EnvSym\then \code{\{$\expr$\}} : \tau
	}
\]

\subsection{Invocation}

\paragraph{Positional}

\[
	\infrule{
		\EnvSym\then n \geq 0 \semiq \forall_{i\in[1,n]}{\tau_i:\Type} \semiq \tau_{n+1}:\Type \\*
		\EnvSym\then \phi = (\tau_1,\dotsc,\tau_n)\to\tau_{n+1} \semiq  f : \phi \\*
		\EnvSym\then \forall_{i\in[1,n]}{\eta_i:\Type,\eta_i\narreq\tau_i,\expr_i:\eta_i}
	}{
		\EnvSym\then \code{$f$($\expr_1$,$\dotsc$,$\expr_n$)} : \tau_{n+1}
	}
\]

\paragraph{Nominal}

\[
	\infrule{
		\EnvSym\then n \geq 0 \semiq \forall_{i\in[1,n]}{x\in\mathcal{N}\subseteq\Id,\tau_i:\Type} \semiq \tau_{n+1}:\Type \\*
		\EnvSym\then \phi = (x_1:\tau_1,\dotsc,x_n:\tau_n)\to\tau_{n+1} \semiq  f : \phi \\*
		\EnvSym\then \forall_{j\in[1,n]}{y\in\mathcal{N},\eta_j:\Type,\expr_j:\eta_j} \\*
		\EnvSym\then \forall_{j\in[1,n]}{\exists_{i\in[1,n]}{y_j=x_i,\eta_j\narreq\tau_i}}
	}{
		\EnvSym\then \code{$f$($y_1$:$\expr_1$,$\dotsc$,$y_n$:$\expr_n$)} : \tau_{n+1}
	}
\]

\comment{TODO: completare}

\subsection{Reference}

\comment{TODO: chiarire notazione e muovere in sezione operatori}

\[
	\infrule{
		\EnvSym\then x \in \Id \\*
		\EnvSym\then n \geq 0 \semiq \forall_{i\in[1,n]}{x_i\in\Id,\tau_i:\Type} \semiq \tau_{n+1}\Type \\*
		\EnvSym\then \phi = (x_1:\tau_1,\dotsc,x_n:\tau_n)\to\tau_{n+1} \\*
	}{
		\EnvSym,\code{$x$ $\phi$}\then \code{::$x$} : \phi
	}
\]

\clearpage

\subsection{Range}

\[
	\infrule{
		\EnvSym\then \tau_1 \in \Number \semiq \expr_1\in\Expr \semiq \expr_1:\tau_1\\*
		\EnvSym\then \tau_2 \in \Number \semiq \expr_2\in\Expr \semiq \expr_2:\tau_2 \\*
		\EnvSym\then \omega = \tau_1 \sqcap \tau_2
	}{
		\EnvSym\then \code{$\expr_1$..$\expr_2$} : \exists\omega.\Range
	}
\]
Range from $\expr_1$ (begin) to $\expr_2$ (end).

\[
	\infrule{
		\EnvSym\then \tau_1 \in \Number \semiq \expr_1\in\Expr \semiq \expr_1:\tau_1\\*
		\EnvSym\then \tau_2 \in \Number \semiq \expr_2\in\Expr \semiq \expr_2:\tau_2 \\*
		\EnvSym\then \tau_3 \in \Number \semiq \expr_3\in\Expr \semiq \expr_3:\tau_3 \\*
		\EnvSym\then \omega = \tau_1 \sqcap \tau_2 \sqcap \tau_3
	}{
		\EnvSym\then \code{$\expr_1$..$\expr_2$..$\expr_3$} : \exists\omega.\Range
	}
\]
Range from $\expr_1$ (begin) to $\expr_2$ (end) with a step $\expr_3$.

\clearpage

\paragraph{General rules}

\[
	\infrule{
		\EnvSym\then \tau_1:\Type \semiq \expr_1:\tau \\*
		\EnvSym\then \tau_2:\Type \semiq \tau_1 \narreq \tau_2 \semiq \expr_2:\tau_2
	}{
		\EnvSym\then \code{$\expr_1$ = $\expr_2$} : \Bool
	}
\]

\[
	\infrule{
		\EnvSym\then \tau_1:\Type \semiq \expr_1:\tau_1 \\*
		\EnvSym\then \tau_2:\Type \semiq \expr_2:\tau_2 \\*
		\EnvSym\then \code{$\expr_1$ = $\expr_2$} : \Bool
	}{
		\EnvSym\then \code{$\expr_1$ != $\expr_2$} : \Bool
	}
\]
If the "equal" operator is defined, the "not equal" operator is also defined.

\vspace{\baselineskip}

\[
	\infrule{
		\EnvSym\then \tau_1:\Type \semiq \expr_1:\tau_1 \\*
		\EnvSym\then \tau_2:\Type \semiq \expr_2:\tau_2 \\*
		\EnvSym\then \code{$\expr_1$ =\~{} $\expr_2$}: \Bool
	}{
		\EnvSym\then \code{$\expr_2$ \~{}= $\expr_1$}: \Bool
	}
\]
If the "contains" operator is defined, the "contained" operator is also
defined.

\vspace{\baselineskip}

\[
	\Number = \set{\Natural, \Int, \Real}
\]

\[
	\infrule{
		\EnvSym\then \expr:\Number \semiq \tau \in \Number
	}{
		\EnvSym\then \code{$\expr$ as $\tau$}: \tau
	}
\]

\comment{TODO(impl): -69 as Natural fallisce (out of bounds). Dovrebbe?}

\[
	\infrule{
		\EnvSym\then \tau \in \Number \semiq \expr:\tau
	}{
		\EnvSym\then \code{+$\expr$}: \tau
	}
\]

\[
	\infrule{
		\EnvSym\then \tau \in \Number\setminus\set{\Natural} \semiq \expr:\tau
	}{
		\EnvSym\then \code{-$\expr$}: \tau
	}
\]

\[
	\infrule{
		\EnvSym\then \tau_1,\tau_2 \in \Number \semiq \expr_1:\tau_1 \semiq \expr_2:\tau_2\\*
		\bop = \set{\code{=}, \code{!=}, \code{<}, \code{<=}, \code{>=}, \code{>=}, \code{>}}
	}{
		\EnvSym\then \code{$\expr_1$ $\bop$ $\expr_2$}: \Bool
	}
\]

\clearpage

\subsection{Primitive types}

\subsubsection{Void}

\comment{TODO: ?}

\subsubsection{Bool}

\[
	\infrule{
		\EnvSym\then \expr:\Bool
	}{
		\EnvSym\then \code{!$\expr$}:\Bool
	}
\]

\[
	\infrule{
		\EnvSym\then \expr_1:\Bool \semiq \expr_2:\Bool \\*
		\bop = \set{\code{and}, \code{\&}, \code{or}, \code{|}}
	}{
		\EnvSym\then \code{$\expr_1$ $\bop$ $\expr_2$}:\Bool
	}
\]

\subsubsection{Natural}

\paragraph{Plus}

\[
	\infrule{
		\EnvSym\then \expr_1:\Natural \semiq \expr_2:\Natural
	}{
		\EnvSym\then \code{$\expr_1$ + $\expr_2$} : \Natural
	}
\]

\[
	\infrule{
		\EnvSym\then \expr_1:\Natural \semiq \expr_2:\Int
	}{
		\EnvSym\then \code{$\expr_1$ + $\expr_2$} : \Int
	}
\]

\[
	\infrule{
		\EnvSym\then \expr_1:\Natural \semiq \expr_2:\Real
	}{
		\EnvSym\then \code{$\expr_1$ + $\expr_2$} : \Real
	}
\]

\paragraph{Minus}

\[
	\infrule{
		\EnvSym\then \expr_1:\Natural \semiq \expr_2:\Natural
	}{
		\EnvSym\then \code{$\expr_1$ - $\expr_2$} : \Int
	}
\]

\comment{TODO: specificare se $l-r\geq0$}

\[
	\infrule{
		\EnvSym\then \expr_1:\Natural \semiq \expr_2:\Int
	}{
		\EnvSym\then \code{$\expr_1$ - $\expr_2$} : \Int
	}
\]

\[
	\infrule{
		\EnvSym\then \expr_1:\Natural \semiq \expr_2:\Real
	}{
		\EnvSym\then \code{$\expr_1$ - $\expr_2$} : \Real
	}
\]

\paragraph{Multiplication}

\[
	\infrule{
		\EnvSym\then \expr_1:\Natural \semiq \expr_2:\Natural
	}{
		\EnvSym\then \code{$\expr_1$ * $\expr_2$} : \Natural
	}
\]


\[
	\infrule{
		\EnvSym\then \expr_1:\Natural \semiq \expr_2:\Int
	}{
		\EnvSym\then \code{$\expr_1$ * $\expr_2$} : \Int
	}
\]

\[
	\infrule{
		\EnvSym\then \expr_1:\Natural \semiq \expr_2:\Real
	}{
		\EnvSym\then \code{$\expr_1$ * $\expr_2$} : \Real
	}
\]

\paragraph{Division}

\[
	\infrule{
		\EnvSym\then \expr_1:\Natural \semiq \expr_2:\Natural
	}{
		\EnvSym\then \code{$\expr_1$ / $\expr_2$} : \Real
	}
\]

\[
	\infrule{
		\EnvSym\then \expr_1:\Natural \semiq \expr_2:\Int
	}{
		\EnvSym\then \code{$\expr_1$ / $\expr_2$} : \Real
	}
\]

\[
	\infrule{
		\EnvSym\then \expr_1:\Natural \semiq \expr_2:\Real
	}{
		\EnvSym\then \code{$\expr_1$ / $\expr_2$} : \Real
	}
\]

\paragraph{Modulo}

\[
	\infrule{
		\EnvSym\then \expr_1:\Natural \semiq \expr_2:\Natural
	}{
		\EnvSym\then \code{$\expr_1$ \% $\expr_2$} : \Natural
	}
\]

\[
	\infrule{
		\EnvSym\then \expr_1:\Natural \semiq \expr_2:\Int
	}{
		\EnvSym\then \code{$\expr_1$ \% $\expr_2$} : \Int
	}
\]

\[
	\infrule{
		\EnvSym\then \expr_1:\Natural \semiq \expr_2:\Real
	}{
		\EnvSym\then \code{$\expr_1$ \% $\expr_2$} : \Real
	}
\]

\paragraph{Power}

\[
	\infrule{
		\EnvSym\then \expr_1:\Natural \semiq \expr_2:\Natural
	}{
		\EnvSym\then \code{$\expr_1$ \^{} $\expr_2$} : \Natural
	}
\]

\[
	\infrule{
		\EnvSym\then \expr_1:\Natural \semiq \expr_2:\Int
	}{
		\EnvSym\then \code{$\expr_1$ \^{} $\expr_2$} : \Int
	}
\]

\[
	\infrule{
		\EnvSym\then \expr_1:\Natural \semiq \expr_2:\Real
	}{
		\EnvSym\then \code{$\expr_1$ \^{} $\expr_2$} : \Real
	}
\]

\subsubsection{Int}

\paragraph{Bitwise not}

\[
	\infrule{
		\EnvSym\then \expr:\Int
	}{
		\EnvSym\then \code{!!$\expr$} : \Int
	}
\]

\paragraph{Bitwise and}

\[
	\infrule{
		\EnvSym\then \expr_1:\Int \semiq \expr_2:\Natural
	}{
		\EnvSym\then \code{$\expr_1$ \&\& $\expr_2$} : \Int
	}
\]

\[
	\infrule{
		\EnvSym\then \expr_1:\Int \semiq \expr_2:\Int
	}{
		\EnvSym\then \code{$\expr_1$ \&\& $\expr_2$} : \Int
	}
\]

\paragraph{Bitwise or}

\[
	\infrule{
		\EnvSym\then \expr_1:\Int \semiq \expr_2:\Natural
	}{
		\EnvSym\then \code{$\expr_1$ || $\expr_2$} : \Int
	}
\]

\[
	\infrule{
		\EnvSym\then \expr_1:\Int \semiq \expr_2:\Int
	}{
		\EnvSym\then \code{$\expr_1$ || $\expr_2$} : \Int
	}
\]

\paragraph{Bitwise xor}

\[
	\infrule{
		\EnvSym\then \expr_1:\Int \semiq \expr_2:\Natural
	}{
		\EnvSym\then \code{$\expr_1$ ||| $\expr_2$} : \Int
	}
\]

\[
	\infrule{
		\EnvSym\then \expr_1:\Int \semiq \expr_2:\Int
	}{
		\EnvSym\then \code{$\expr_1$ ||| $\expr_2$} : \Int
	}
\]

\paragraph{Bitwise left shift}

\[
	\infrule{
		\EnvSym\then \expr_1:\Int \semiq \expr_2:\Natural
	}{
		\EnvSym\then \code{$\expr_1$ \{<\} $\expr_2$} : \Int
	}
\]

\[
	\infrule{
		\EnvSym\then \expr_1:\Int \semiq \expr_2:\Int
	}{
		\EnvSym\then \code{$\expr_1$ \{<\} $\expr_2$} : \Int
	}
\]

\paragraph{Bitwise right shift}

\[
	\infrule{
		\EnvSym\then \expr_1:\Int \semiq \expr_2:\Natural
	}{
		\EnvSym\then \code{$\expr_1$ \{>\} $\expr_2$} : \Int
	}
\]

\[
	\infrule{
		\EnvSym\then \expr_1:\Int \semiq \expr_2:\Int
	}{
		\EnvSym\then \code{$\expr_1$ \{>\} $\expr_2$} : \Int
	}
\]

\paragraph{Bitwise right unsigned shift}

\[
	\infrule{
		\EnvSym\then \expr_1:\Int \semiq \expr_2:\Natural
	}{
		\EnvSym\then \code{$\expr_1$ \{>>\} $\expr_2$} : \Int
	}
\]

\[
	\infrule{
		\EnvSym\then \expr_1:\Int \semiq \expr_2:\Int
	}{
		\EnvSym\then \code{$\expr_1$ \{>>\} $\expr_2$} : \Int
	}
\]

\paragraph{Plus}

\[
	\infrule{
		\EnvSym\then \expr_1:\Int \semiq \expr_2:\Natural
	}{
		\EnvSym\then \code{$\expr_1$ + $\expr_2$} : \Int
	}
\]

\[
	\infrule{
		\EnvSym\then \expr_1:\Int \semiq \expr_2:\Int
	}{
		\EnvSym\then \code{$\expr_1$ + $\expr_2$} : \Int
	}
\]

\[
	\infrule{
		\EnvSym\then \expr_1:\Int \semiq \expr_2:\Real
	}{
		\EnvSym\then \code{$\expr_1$ + $\expr_2$} : \Real
	}
\]

\paragraph{Minus}

\[
	\infrule{
		\EnvSym\then \expr_1:\Int \semiq \expr_2:\Natural
	}{
		\EnvSym\then \code{$\expr_1$ - $\expr_2$} : \Int
	}
\]

\[
	\infrule{
		\EnvSym\then \expr_1:\Int \semiq \expr_2:\Int
	}{
		\EnvSym\then \code{$\expr_1$ - $\expr_2$} : \Int
	}
\]

\[
	\infrule{
		\EnvSym\then \expr_1:\Int \semiq \expr_2:\Real
	}{
		\EnvSym\then \code{$\expr_1$ - $\expr_2$} : \Real
	}
\]

\paragraph{Multiplication}

\[
	\infrule{
		\EnvSym\then \expr_1:\Int \semiq \expr_2:\Natural
	}{
		\EnvSym\then \code{$\expr_1$ * $\expr_2$} : \Int
	}
\]

\[
	\infrule{
		\EnvSym\then \expr_1:\Int \semiq \expr_2:\Int
	}{
		\EnvSym\then \code{$\expr_1$ * $\expr_2$} : \Int
	}
\]

\[
	\infrule{
		\EnvSym\then \expr_1:\Int \semiq \expr_2:\Real
	}{
		\EnvSym\then \code{$\expr_1$ * $\expr_2$} : \Real
	}
\]

\paragraph{Division}

\[
	\infrule{
		\EnvSym\then \expr_1:\Int \semiq \expr_2:\Natural
	}{
		\EnvSym\then \code{$\expr_1$ / $\expr_2$} : \Real
	}
\]

\[
	\infrule{
		\EnvSym\then \expr_1:\Int \semiq \expr_2:\Int
	}{
		\EnvSym\then \code{$\expr_1$ / $\expr_2$} : \Real
	}
\]

\[
	\infrule{
		\EnvSym\then \expr_1:\Int \semiq \expr_2:\Real
	}{
		\EnvSym\then \code{$\expr_1$ / $\expr_2$} : \Real
	}
\]

\paragraph{Modulo}

\[
	\infrule{
		\EnvSym\then \expr_1:\Int \semiq \expr_2:\Natural
	}{
		\EnvSym\then \code{$\expr_1$ \% $\expr_2$} : \Int
	}
\]

\[
	\infrule{
		\EnvSym\then \expr_1:\Int \semiq \expr_2:\Int
	}{
		\EnvSym\then \code{$\expr_1$ \% $\expr_2$} : \Int
	}
\]

\[
	\infrule{
		\EnvSym\then \expr_1:\Int \semiq \expr_2:\Real
	}{
		\EnvSym\then \code{$\expr_1$ \% $\expr_2$} : \Real
	}
\]

\paragraph{Power}

\[
	\infrule{
		\EnvSym\then \expr_1:\Int \semiq \expr_2:\Natural
	}{
		\EnvSym\then \code{$\expr_1$ \^{} $\expr_2$} : \Real
	}
\]

\[
	\infrule{
		\EnvSym\then \expr_1:\Int \semiq \expr_2:\Int
	}{
		\EnvSym\then \code{$\expr_1$ \^{} $\expr_2$} : \Real
	}
\]

\[
	\infrule{
		\EnvSym\then \expr_1:\Int \semiq \expr_2:\Real
	}{
		\EnvSym\then \code{$\expr_1$ \^{} $\expr_2$} : \Real
	}
\]

\subsubsection{Real}

\paragraph{Plus}

\[
	\infrule{
		\EnvSym\then \expr_1:\Int \semiq \tau \in \Number \semiq \expr_2:\tau
	}{
		\EnvSym\then \code{$\expr_1$ + $\expr_2$} : \Real
	}
\]

\paragraph{Minus}

\[
	\infrule{
		\EnvSym\then \expr_1:\Int \semiq \tau \in \Number \semiq \expr_2:\tau
	}{
		\EnvSym\then \code{$\expr_1$ - $\expr_2$} : \Real
	}
\]

\paragraph{Multiplication}

\[
	\infrule{
		\EnvSym\then \expr_1:\Int \semiq \tau \in \Number \semiq \expr_2:\tau
	}{
		\EnvSym\then \code{$\expr_1$ * $\expr_2$} : \Real
	}
\]

\paragraph{Division}

\[
	\infrule{
		\EnvSym\then \expr_1:\Int \semiq \tau \in \Number \semiq \expr_2:\tau
	}{
		\EnvSym\then \code{$\expr_1$ / $\expr_2$} : \Real
	}
\]

\paragraph{Modulo}

\[
	\infrule{
		\EnvSym\then \expr_1:\Int \semiq \tau \in \Number \semiq \expr_2:\tau
	}{
		\EnvSym\then \code{$\expr_1$ \% $\expr_2$} : \Real
	}
\]

\paragraph{Power}

\[
	\infrule{
		\EnvSym\then \expr_1:\Int \semiq \tau \in \Number \semiq \expr_2:\tau
	}{
		\EnvSym\then \code{$\expr_1$ \^{} $\expr_2$} : \Real
	}
\]

\subsubsection{Text}

\paragraph{As}

\[
	\infrule{
		\EnvSym\then \expr_1:\Text \semiq \tau \in \Number \cup \set{\Text, \Bool}
	}{
		\EnvSym\then \code{$\expr_1$ as $\tau$}: \tau
	}
\]

\paragraph{Plus}

\[
	\infrule{
		\EnvSym\then \expr_1:\Text \semiq \tau:\Type \semiq \expr_2:\tau
	}{
		\EnvSym\then \code{$\expr_1$ + $\expr_2$}: \Text
	}
\]

\paragraph{Minus}

\[
	\infrule{
		\EnvSym\then \expr_1:\Text \semiq \expr_2:\Text
	}{
		\EnvSym\then \code{$\expr_1$ - $\expr_2$}: \Text
	}
\]

\paragraph{Power}

\[
	\infrule{
		\EnvSym\then \expr_1:\Text \semiq \tau \in \Number \semiq \expr_2:\tau
	}{
		\EnvSym\then \code{$\expr_1$ \^{} $\expr_2$}: \Text
	}
\]

\comment{TODO(impl): non dovremmo permettere Real}

\paragraph{Division}

\[
	\infrule{
		\EnvSym\then \expr_1:\Text \semiq \expr_2:\Text
	}{
		\EnvSym\then \code{$\expr_1$ / $\expr_2$}: \exists\Text.\List
	}
\]

\paragraph{Equal}

\[
	\infrule{
		\EnvSym\then \expr_1:\Text \semiq \expr_2:\Text
	}{
		\EnvSym\then \code{$\expr_1$ = $\expr_2$}: \Bool
	}
\]

\paragraph{Less}

\[
	\infrule{
		\EnvSym\then \expr_1:\Text \semiq \expr_2:\Text
	}{
		\EnvSym\then \code{$\expr_1$ < $\expr_2$}: \Bool
	}
\]

\paragraph{Less equal}

\[
	\infrule{
		\EnvSym\then \expr_1:\Text \semiq \expr_2:\Text
	}{
		\EnvSym\then \code{$\expr_1$ <= $\expr_2$}: \Bool
	}
\]

\paragraph{Greater equal}

\[
	\infrule{
		\EnvSym\then \expr_1:\Text \semiq \expr_2:\Text
	}{
		\EnvSym\then \code{$\expr_1$ >= $\expr_2$}: \Bool
	}
\]

\paragraph{Greater}

\[
	\infrule{
		\EnvSym\then \expr_1:\Text \semiq \expr_2:\Text
	}{
		\EnvSym\then \code{$\expr_1$ > $\expr_2$}: \Bool
	}
\]

\paragraph{Contains}

\[
	\infrule{
		\EnvSym\then \expr_1:\Text \semiq \expr_2:\Text
	}{
		\EnvSym\then \code{$\expr_1$ =\~{} $\expr_2$}: \Bool
	}
\]

\paragraph{Dot access}

\[
	\infrule{
		\text{Index access}\\*
		\EnvSym\then \expr_1:\Text \semiq \expr_2:\Int
	}{
		\EnvSym\then \code{$\expr_1$.$\expr_2$}: \Text
	}
\]

\[
	\infrule{
		\text{"Slice"}\\*
		\EnvSym\then \expr_1:\Text \semiq \expr_2:\exists\Int.\Range
	}{
		\EnvSym\then \code{$\expr_1$.$\expr_2$}: \Text
	}
\]

\[
	\infrule{
		\text{Text length}\\*
		\EnvSym\then \expr_1:\Text
	}{
		\EnvSym\then \code{$\expr_1$.\_count}: \Int
	}
\]

\[
	\infrule{
		\text{Characters range}\\*
		\EnvSym\then \expr_1:\Text
	}{
		\EnvSym\then \code{$\expr_1$.\_range}: \exists\Int.\Range
	}
\]

\[
	\infrule{
		\text{Characters list}\\*
		\EnvSym\then \expr_1:\Text
	}{
		\EnvSym\then \code{$\expr_1$.\_chars}: \exists\Text.\List
	}
\]

\clearpage

\subsection{Collection types}

\subsubsection{List}

\paragraph{As}

\[
	\infrule{
		\EnvSym\then \tau_1 = \List[\nu_1 := \Type] \\*
		\EnvSym\then \tau_2 = \List[\nu_2 := \Type] \\*
		\EnvSym\then \expr:\tau_1 \semiq x:\nu_1 \\*
		\EnvSym\then \code{$x$ as $\nu_2$} : \nu_2
	}{
		\EnvSym\then \code{$\expr$ as $\tau_2$} : \tau_2
	}
\]
If a term of type $\nu_1$ can be casted to type $\nu_2$ $\implies$ A list
of $\nu_1$ can be casted to a list of $\nu_2$.

\paragraph{Plus}

\[
	\infrule{
		\EnvSym\then \tau_1=\List \semiq \expr_1:\tau_1 \\*
		\EnvSym\then \tau_2=\List \semiq \expr_2:\tau_2 \\*
		\EnvSym\then \tau_2 \narreq \tau_1
	}{
		\EnvSym\then \code{$\expr_1$ + $\expr_2$} : \tau_1
	}
\]
\comment{TODO: chiarire. Il tipo della concatenazione dovrebbe essere deciso da
lhs o dovrebbe essere il widest tra i due?}

\[
	\infrule{
		\EnvSym\then \tau=\List[\nu_1:=\Type] \semiq \expr_1:\tau \\*
		\EnvSym\then \nu_2:\Type \semiq \nu_2\narreq\nu_1 \semiq \expr_2:\nu_2
	}{
		\EnvSym\then \code{$\expr_1$ + $\expr_2$} : \tau
	}
\]

\paragraph{Minus}

\[
	\infrule{
		\EnvSym\then \tau_1=\List \semiq \expr_1:\tau_1 \\*
		\EnvSym\then \tau_2=\List \semiq \expr_2:\tau_2 \\*
		\EnvSym\then \tau_2 \narreq \tau_1
	}{
		\EnvSym\then \code{$\expr_1$ - $\expr_2$} : \tau_1
	}
\]
\comment{TODO: chiarire. Il tipo dell'intersezione dovrebbe essere deciso da
lhs o dovrebbe essere il widest tra i due?}

\[
	\infrule{
		\EnvSym\then \tau=\List[\nu_1:=\Type] \semiq \expr_1:\tau \\*
		\EnvSym\then \nu_2:\Type \semiq \nu_2\narreq\nu_1 \semiq \expr_2:\nu_2
	}{
		\EnvSym\then \code{$\expr_1$ - $\expr_2$} : \tau
	}
\]

\paragraph{Contains}

\[
	\infrule{
		\EnvSym\then \tau_1:\Type \semiq \expr_1:\exists\tau_1.\List \\*
		\EnvSym\then \tau_2:\Type \semiq \tau_2 \narreq \tau_1 \semiq \expr_2:\tau_2
	}{
		\EnvSym\then \code{$\expr_1$ =\~{} $\expr_2$} : \Bool
	}
\]
A list of $\tau_1$ contains an element of type $\tau_2$ (narrower than
$\tau_1$).

\[
	\infrule{
		\EnvSym\then \tau_1:\Type \semiq \expr_1:\exists\tau_1.\List\\*
		\EnvSym\then \tau_2:\Type \semiq \tau_2 \narreq \tau_1 \semiq \expr_2:\exists\tau_2.\List
	}{
		\EnvSym\then \code{$\expr_1$ =\~{} $\expr_2$} : \Bool
	}
\]
A list of $\tau_1$ contains the elements of a list of $\tau_2$ (narrower than
$\tau_1$).

\paragraph{Dot access}

\[
	\infrule{
		\text{Index access}\\*
		\EnvSym\then \tau:\Type
		\semiq \expr_1:\exists\tau.\List
		\semiq \expr_2:\Int
	}{
		\EnvSym\then \code{$\expr_1$.$\expr_2$}: \tau
	}
\]

\[
	\infrule{
		\text{"Slice"}\\*
		\EnvSym\then \tau : \Type
		\semiq \expr_1:\exists\tau.\List
		\semiq \expr_2:\exists\Int.\Range
	}{
		\EnvSym\then \code{$\expr_1$.$\expr_2$}: \tau
	}
\]

\[
	\infrule{
		\text{Count access}\\*
		\EnvSym\then \tau:\Type \semiq \expr:\exists\tau.\List
	}{
		\EnvSym\then \code{$\expr$.\_count}: \Int
	}
\]

\[
	\infrule{
		\text{Range access}\\*
		\EnvSym\then \expr:\Text
	}{
		\EnvSym\then \code{$\expr$.\_range}: \exists\Int.\Range
	}
\]

\subsubsection{Map}

\paragraph{As}

\[
	\infrule{
		\EnvSym\then \tau_1 = \Map[\kappa_1:=\Type,\nu_1:=\Type] \\*
		\EnvSym\then \tau_2 = \Map[\kappa_2:=\Type,\nu_2:=\Type] \\*
		\EnvSym\then \expr:\tau_1 \semiq \expr_k:\kappa_1 \semiq \expr_v:\nu_1 \\*
		\EnvSym\then \code{$\expr_k$ as $\kappa_2$} : \kappa_2 \land \code{$\expr_v$ as $\nu_2$} : \nu_2
	}{
		\EnvSym\then \code{$\expr$ as $\tau_2$} : \tau_2
	}
\]
If a term of type $\kappa_1$ and be casted to type $\kappa_2$ $\land$ A term of
type $\nu_1$ can be casted to type $\nu_2$ $\implies$ A map of type $\tau_1$
can be casted to the type $\tau_2$.

\paragraph{Plus}

\[
	\infrule{
		\EnvSym\then \tau_1=\Map \semiq \expr_1:\tau_1 \\*
		\EnvSym\then \tau_2=\Map \semiq \expr_2:\tau_2 \\*
		\EnvSym\then \tau_2 \narreq \tau_1
	}{
		\EnvSym\then \code{$\expr_1$ + $\expr_2$} : \tau_1
	}
\]
\comment{TODO: chiarire. Il tipo della concatenazione dovrebbe essere deciso da
lhs o dovrebbe essere il widest tra i due?}

\paragraph{Minus}

\[
	\infrule{
		\EnvSym\then \tau_1=\Map \semiq \expr_1:\tau_1 \\*
		\EnvSym\then \tau_2=\Map \semiq \expr_2:\tau_2 \\*
		\EnvSym\then \tau_2 \narreq \tau_1
	}{
		\EnvSym\then \code{$\expr_1$ - $\expr_2$} : \tau_1
	}
\]
\comment{TODO: chiarire. Il tipo dell'intersezione dovrebbe essere deciso da
lhs o dovrebbe essere il widest tra i due?}

\[
	\infrule{
		\EnvSym\then \tau=\Map[\kappa_1:=\Type] \semiq \expr_1:\tau \\*
		\EnvSym\then \kappa_2:\Type \semiq \kappa_2\narreq\kappa_1 \semiq \expr_2:\kappa_2
	}{
		\EnvSym\then \code{$\expr_1$ - $\expr_2$} : \tau
	}
\]


\paragraph{Contains}

\[
	\infrule{
		\EnvSym\then \tau = \Map[\kappa_1:=\Type,\nu:=\Type] \semiq \expr_1:\tau \\*
		\EnvSym\then \kappa_2:\Type \semiq \kappa_2 \narreq \kappa_1 \semiq \expr_2:\kappa_2
	}{
		\EnvSym\then \code{$\expr_1$ =\~{} $\expr_2$} : \Bool
	}
\]
A map $\expr_1$ of type $\tau$ contains a key of type $\kappa_2$ (narrower than
$\kappa_1$).

\[
	\infrule{
		\EnvSym\then \tau_1 = \Map[\kappa_1:=\Type,\nu_1:=\Type] \semiq \expr_1:\tau_1 \\*
		\EnvSym\then \tau_2 = \Map[\kappa_2:=\Type,\nu_2:=\Type] \semiq \expr_2:\tau_2 \\*
		\EnvSym\then \kappa_2 \narreq \kappa_1 \land \nu_2 \narreq \nu_1
	}{
		\EnvSym\then \code{$\expr_1$ =\~{} $\expr_2$} : \Bool
	}
\]
A map $\expr_1$ of type $\tau_1$ contais the key-value pairs of type $\kappa_2$
(narrower than $\kappa_1$) and $\nu_2$ (narrower than $\nu_1$) of the map
$\expr_2$ of type $\tau_2$.

\paragraph{Dot access}

\[
	\infrule{
		\text{Access by key}\\*
		\EnvSym\then \tau = \Map[\kappa_1:=\Type,\nu:=\Type] \semiq \expr_1:\tau \\*
		\EnvSym\then \kappa_2 = \Type \semiq \kappa_2 \narreq \kappa_1 \semiq \expr_2:\kappa_2
	}{
		\EnvSym\then \code{$\expr_1$.$\expr_2$}: \nu
	}
\]

\[
	\infrule{
		\text{Count access}\\*
		\EnvSym\then \tau = \Map[\kappa:=\Type,\nu:=\Type] \semiq \expr:\tau \\*
	}{
		\EnvSym\then \code{$\expr$.\_count}: \Int
	}
\]

\[
	\infrule{
		\text{Keys access}\\*
		\EnvSym\then \tau = \Map[\kappa:=\Type,\nu:=\Type] \semiq \expr:\tau \\*
	}{
		\EnvSym\then \code{$\expr$.\_keys}: \exists\kappa.\List
	}
\]

\[
	\infrule{
		\text{Values access}\\*
		\EnvSym\then \tau = \Map[\kappa:=\Type,\nu:=\Type] \semiq \expr:\tau \\*
	}{
		\EnvSym\then \code{$\expr$.\_values}: \exists\nu.\List
	}
\]

\subsection{Product types}

\[
	\Product = \set{\texttt{Tuple, Record, Struct}}
\]

\[
	\infrule{
		\EnvSym\then \tau_1 \in \Product \semiq \expr_1:\tau_1 \\*
		\EnvSym\then \tau_2 \in \Product \semiq \expr_2:\tau_2
	}{
		\EnvSym\then \code{$\expr_1$ = $\expr_2$}: \Bool
	}
\]

\[
	\infrule{
		\text{Fields count access} \\*
		\EnvSym\then \tau\in\Product \semiq \expr:\tau
	}{
		\EnvSym\then \code{$\expr$.\_count}: \Int
	}
\]

\[
	\infrule{
		\text{Fields range access} \\*
		\EnvSym\then \tau\in\Product \semiq \expr:\tau
	}{
		\EnvSym\then \code{$\expr$.\_range}: \exists\Int.\Range
	}
\]

\[
	\infrule{
		\text{Fields names access} \\*
		\EnvSym\then \tau\in\set{\Record,\Struct} \semiq \expr:\tau
	}{
		\EnvSym\then \code{$\expr$.\_fields}: \exists\Text.\List
	}
\]

\subsubsection{Tuple}

\paragraph{Dot access}

\[
	\infrule{
		\text{"Slice"}\\*
		\EnvSym\then n \geq 0 \semiq \forall_{i\in[1,n]}{\tau_i:\Type} \semiq \expr_1:\tau_1 \times \dotsb \times \tau_n \\*
		\EnvSym\then \expr_2:\exists\Int.\Range
	}{
		\EnvSym\then \code{$\expr_1$.$\expr_2$}: \bigtimes_{i \in \expr_2} \tau_i
	}
\]

Example: \code{<0,1,2,3,4,5> . 2..4 = <2 3 4>}

With step: \code{<0,1,2,3,4,5> . 0..5..2 = <0 2 4>}

Reverse: \code{<0,1,2,3,4,5> . 4..2 = <4 3 2>}

\[
	\infrule{
		\text{Access by index}\\*
		\EnvSym\then n \geq 0 \semiq \forall_{i\in[1,n]}{\tau_i:\Type} \semiq \expr_1:\tau_1 \times \dotsb \times \tau_n \\*
		\EnvSym\then \expr_2:\Int \semiq \expr_2 \in [1,n]
	}{
		\EnvSym\then \code{$\expr_1$.$\expr_2$}: \tau_{\expr_2}
	}
\]

\subsubsection{Record}

\paragraph{Dot access}

\[
	\infrule{
		\text{Access by index} \\*
		\EnvSym\then n \geq 1 \semiq \forall_{i\in[1,n]}{\tau_i:\Type, x_i\in\Id} \semiq \expr_1:\tuple{x_1:\tau_1,\dotsc,x_n:\tau_n}\\*
		\EnvSym\then \expr_2:\Int \semiq \expr_2 \in [1,n]
	}{
		\EnvSym\then \code{$\expr_1$.$\expr_2$}: \tau_{\expr_2}
	}
\]

\comment{TODO(impl): permettere solo a Tuple in modo da rendere esplicito la perdita degli identificatori}

\[
	\infrule{
		\text{"Slice"}\\*
		\EnvSym\then n \geq 1 \semiq \forall_{i\in[1,n]}{\tau_i:\Type, x_i\in\Id} \semiq \rho=\tuple{x_1:\tau_1,\dotsc,x_n:\tau_n} \semiq \expr_1:\rho\\*
		\EnvSym\then \expr_2:\exists\Int.\Range \\*
		\EnvSym\then \forall_{i\in\expr_2}{\rho'=\tuple{x_j:\tau_j, \dotsc, x_j:\tau_j}}
	}{
		\EnvSym\then \code{$\expr_1$.$\expr_2$}: \rho'
	}
\]

\[
	\infrule{
		\text{Access by name}\\*
		\EnvSym\then n \geq 1 \semiq \forall_{i\in[1,n]}{\tau_i:\Type, x_i \in \Id} \semiq \expr:\tuple{x_1:\tau_1,\dotsc,x_n:\tau_n} \\*
		\EnvSym\then x' \in \Id \semiq \exists_{j\in[1,n]} x_j = x'
	}{
		\EnvSym\then \code{$\expr$.$x'$}: \tau_j
	}
\]

\subsubsection{Struct}

\comment{TODO: ...}

\subsection{Function}

\comment{TODO: ...}

\clearpage

\section{Statements}

\subsection{Variable declaration}

\[
	\infrule{
		\EnvSym\then x \in \Id \semiq \tau:\Type
	}{
		\EnvSym, \code{$x$ $\tau$} \then x:\tau
	}
\]

\[
	\infrule{
		\EnvSym\then x\in\Id \semiq \tau:\Type \\*
		\EnvSym\then \tau_\expr:\Type \semiq \tau_\expr\narreq\tau \semiq \expr\in\Expr \semiq \expr:\tau_\expr
	}{
		\EnvSym, \code{var $x$ $\tau$ <- $\expr$} \then x:\tau \qquad \text{(esplicit type)} \\*
		\land \EnvSym, \code{var $x$ <- $\expr$} \then x:\tau \qquad \text{(implicit type)}
	}
\]

\comment{TODO: specificare il tipo Var come da implementazione?}

\subsection{Constant declaration}

\[
	\infrule{
		\EnvSym\then x\in\Id \semiq \tau:\Type \\*
		\EnvSym\then \tau_\expr:\Type \semiq \tau_\expr\narreq\tau \semiq \expr\in\Expr \semiq \expr:\tau_\expr
	}{
		\EnvSym, \code{const $x$ $\tau$ <- $\expr$} \then x:\tau \qquad \text{(esplicit type)} \\*
		\land \EnvSym, \code{const $x$ <- $\expr$} \then x:\tau \qquad \text{(implicit type)}
	}
\]

\subsection{Alias declaration}

\[
	\infrule{
		\EnvSym\then x \in \Id \semiq \expr \in \Expr \semiq \tau:\Type \semiq \expr:\tau
	}{
		\EnvSym, \code{alias $e$ as $x$} \then x:\tau
	}
\]

\subsection{For loop}

\[
	\infrule{
		\EnvSym\then x \in \Id \semiq \tau:\Type \semiq \expr:\exists\tau.\List
	}{
		\EnvSym, \code{for $x$ in $\expr$ do} \then x:\tau
	}
\]

\[
	\infrule{
		\EnvSym\then k \in \Id \semiq v \in \Id \semiq \kappa:\Type \semiq \nu:\Type \semiq \expr:\exists\kappa.\exists\nu.\Map
	}{
		\EnvSym, \code{for $k$=>$v$ in $\expr$ do} \then k:\kappa \semiq v:\nu
	}
\]

\comment{TODO: deconstruction}

\chapter{Operational semantics}
\label{ch:opsem}

This section provides a semi-formal definition of the operational semantics of
the language, definining the state and context and how they mutate with every
expression or statement.

\section{Notation and terminology}

\paragraph{Common sets} Some sets will be used throughout the rules
\[
	\begin{split}
		\Expr &= \set{\expr \mid \expr := \text{any expression}} \\*
		\Val  &= \set{\val \mid \val := \text{any value}} \\*
	\end{split}
\]

\paragraph{Expression} The $\expr$ will be used to denote any expression
permitted by the language (e.g. a literal, a binary expression, a function
invocation etc.).\\*
Unless noted otherwise, $\expr \in \Expr$ is implied.

\paragraph{Expression's value} To denote to what value an expression evaluates
to, the $\eval$ symbol will be used, for example: $\expr \eval \val$\\*
To specify the value along with its type, the notation $\tau(\val)$ will be
used, where $\tau:\Type$ and $\val\in\Val$, for example:
$\expr\eval\Int(\val)$\\*
Unless noted otherwise, $\val \in \Val$ is implied.

\paragraph{Environments} Two environments will be used:
\begin{itemize}
	\item $\EnvLoc :=$ \textbf{location environment}. Describes the
		mappings/bindings between identifiers or an evaluated expression and
		their "location/address" in memory.
		\[
			\begin{split}
				\EnvLoc: \Id \cup \Expr & \to \mathbb{Z}\\*
				\EnvLoc(x) & \mapsto \addr
			\end{split}
		\]
		An address will usually be denoted by $\addr$
	\item $\EnvStore :=$ \textbf{store environment}. Describes the mappings
		between a location in memory and its data content
		\[
			\begin{split}
				\EnvStore : \mathbb{Z} & \to \Val \\*
				\EnvStore(\addr) & \mapsto \val
			\end{split}
		\]
\end{itemize}

\nb{
	"Location" and "memory" are used loosely here. They do not represent an
	actual computer memory, but rather an abstraction of it.
}

\paragraph{Expression's location} The expression location and its value
mappings in the environment is implied
\[
	\expr\eval\val
	\implies \exists_\addr{\EnvLoc(\expr)=\addr \land \EnvStore(\addr)=\val}
\]

\paragraph{Values identity} In this context, values do not have an identity,
meaning that if two memory locations are mapped to the same value, the value
data is effectively duplicated.

\paragraph{Unmapped values} A slight abuse of notation is used to denote the
lack of mapping in an environment via the $\varnothing$ symbol.\\*
For example: $\EnvLoc(i) = \varnothing, \EnvStore(\addr) = \varnothing$

\paragraph{Environment extension} Like in type inference rules, to denote an
environment extending, modifying or overriding an existing one, this notation
will be used $\EnvLoc\env{i\mapsto\addr}$ where it's defined as follows:
\[
	\EnvLoc\env{i\mapsto\addr}(j) \mapsto \begin{cases}
		\addr      &\text{ if } i = j\\
		\EnvLoc(j) &\text{ if } i \neq j\\
	\end{cases}
\]

\paragraph{The $\alloc()$ function} When store environment is extended with
a mapping, it's implicit that it is added if it does not exist, and it is
updated if it already exists. In some cases, to improve clarity, the $\alloc()$
function will be used to explicitly state that a new memory "address" is
allocated in the specified environment. The $\alloc()$ function takes a store
environment and returns the newly allocated address.\\*
For example: $\addr = \alloc(\EnvStore)$

% \subparagraph{The $\copi()$ function} The language (almost) never pass values
% by reference, hence many copies are made during an execution. The simplify the
% rules, a $\copi()$ function is used as a shorthand for allocating a new address
% and mapping it a copy of the specified value

\paragraph{Environment modification} If the evaluation of an expression
results in an alteration of an environment, the judgment may be followed by a
"double turnstile" symbol $\Then$ and the new environments.\\*
For example: $\EnvLoc,\EnvStore \then \expr \Then \EnvStore_2$

\clearpage

\subsection{Zero values}
\label{sec:zero_values}

For more info, see \See{ch:defaults}.

A type's zero value will be described using $\Zero$
\[
	\begin{split}
		\Zero : \Type &\to \Val \\*
		\Zero(\tau)   &\mapsto \val
	\end{split}
\]

Since in some cases the zero value may be overridden, the $\Zero$ function can
be applied on an identifier
\[
	\begin{split}
		\Zero : \Id &\to \Val \\*
		\Zero(i) &\mapsto \begin{cases}
			\val             &\text{ if default overridden } \\*
			\Zero(\EnvSym(i)) &
		\end{cases}
	\end{split}
\]

\subsubsection{Primitive types}

\[
	\begin{split}
		\EnvSym\then \Zero(\Bool)    &= \Bool(\texttt{false}) \\*
		\EnvSym\then \Zero(\Text)    &= \Text(\texttt{""}) \\*
		\EnvSym\then \Zero(\Natural) &= \Natural(\texttt{0}) \\*
		\EnvSym\then \Zero(\Int)     &= \Int(\texttt{0}) \\*
		\EnvSym\then \Zero(\Real)    &= \Real(\texttt{0}) \\*
		\EnvSym\then \Zero(\Type)    &= \Type(\Void) \\*
		\EnvSym\then \Zero(\Void)    &= \Void() \\*
	\end{split}
\]

\subsubsection{Collection types}

\paragraph{List}

\[
	\EnvSym\then \Zero(\List[\nu:=\Type]) = \List(0, \Zero(\nu))
\]

\comment{TODO: se eventualmente formalizzo l'operatore \code{**}, aggiornare
dimensione iniziale}

\paragraph{Map}

\[
	\EnvSym\then \Zero(\Map[\kappa:=\Type,\nu:=\Type]) = \Map(\Zero(\kappa), \Zero(\nu))
\]

\subsubsection{Product types}

\comment{TODO: riscrivere usando environment per membri (con default override)}

\paragraph{Tuple}

\[
	\infrule{
		\EnvSym\then \forall_{i\in[1,n]}{\tau_i:\Type} \semiq n \geq 0
	}{
		\EnvSym\then \Zero(\tau_1 \times \dotsb \times \tau_n) = \Zero(\tau_1) \times \dotsb \times \Zero(\tau_n)
	}
\]

\paragraph{Record}

\[
	\infrule{
		\EnvSym\then n \geq 0 \semiq \forall_{i\in[1,n]}{\tau_i:\Type, x_i \in \Id}
	}{
		\EnvSym\then \Zero(\tuple{x_1:\tau_1,\dotsc,x_n:\tau_n}) = \tuple{x_1:\Zero(\tau_1), \dotsc, x_n:\Zero(\tau_n)}
	}
\]

\paragraph{Struct}

\text{}

\comment{TODO: completare}

\subsubsection{Function}

\[
	\EnvSym\then \Zero(\Function) = ()\to\Void
\]

% \[
% 	\infrule{
% 		\Env\then n \in \mathbb{N} \semiq \forall_{i\in[1,n]}{\tau_i:\Type,x_i\in\Id} \semiq \rho:\Type \\*
% 		\Env\then \phi : (\tau_1,\dotsc,\tau_n)\to\rho
% 	}{
% 		\Env\then \Zero(\phi()) = \Zero(\rho)
% 	}
% \]

\subsubsection{Range}

\[
	\EnvSym\then \Zero(\Range) = \Range(0, 0, 1)
\]

\clearpage

\section{Operational rules}

\subsection{Literals}

\comment{TODO: come nelle regole di inferenza?}

\paragraph{Bool}

\[
	\oprule{
		b \in \set{\texttt{true}, \texttt{false}}
	}{
		\EnvLoc,\EnvStore\then \code{$b$} \eval \texttt{Bool($b$)}
	}
\]

\paragraph{Natural}

\[
	\oprule{
		n := \text{ natural literal }
	}{
		\EnvLoc,\EnvStore\then \code{$n$} \eval \texttt{Natural($n$)}
	}
\]

\paragraph{Int}

\[
	\oprule{
		i := \text{ integer literal }
	}{
		\EnvLoc,\EnvStore\then \code{$i$} \eval \texttt{Int($i$)}
	}
\]

\paragraph{Real}

\[
	\oprule{
		r := \text{ real literal }
	}{
		\EnvLoc,\EnvStore\then \code{$r$} \eval \texttt{Real($r$)}
	}
\]


\paragraph{Text}

\[
	\oprule{
		t := \text{ text literal }
	}{
		\EnvLoc,\EnvStore\then \code{$t$} \eval \texttt{Text($t$)}
	}
\]

\comment{TODO: variables interpolation?}

\subsection{Identifiers}

\[
	\oprule{
		x \in \Id \semiq \EnvLoc(x) = \addr \semiq \EnvStore(\addr) = \val
	}{
		\EnvLoc,\EnvStore\then \code{$x$} \eval \val
	}
\]

\subsection{Statements}

\paragraph{Redirection (assignment)}

\[
	\oprule[aligned]{
		\EnvLoc,\EnvStore     &\then x \in \Id \semiq \EnvLoc(x) = \addr_x \semiq \EnvStore(\addr_x) = \val_x \\*
		\EnvLoc,\EnvStore     &\then \expr \eval \val \Then \EnvLoc_2,\EnvStore_2 \\*
		\EnvLoc_2,\EnvStore_2 &\then \addr_x' = \alloc(\EnvStore_2) \\*
	}{
		\EnvLoc_2,\EnvStore_2\then \code{$x$ <- $\expr$} \Then \EnvLoc_2\env{x \mapsto \addr_x'},\EnvStore_2\env{\addr_x' \mapsto \val}
	}
\]

\paragraph{Variable declaration}

\[
	\oprule[aligned]{
		\EnvSym           &\then x \in \Id \semiq \tau : \Type \\*
		\EnvLoc,\EnvStore &\then \EnvLoc(x) = \varnothing \\*
		\EnvLoc,\EnvStore &\then \addr = \alloc(\EnvStore)
	}{
		\EnvSym,\EnvLoc,\EnvStore \then \code{$x$ $\tau$}
		\Then \EnvLoc\env{x \mapsto \addr},\EnvStore\env{\addr \mapsto \Zero(\tau)}
	}
\]

\subparagraph{Initialization} A variable can be initialized at declaration
\[
	\oprule[aligned]{
		\EnvSym               &\then x \in \Id \semiq \tau : \Type \\*
		\EnvLoc,\EnvStore     &\then \EnvLoc(x) = \varnothing \\*
		\EnvLoc,\EnvStore     &\then \expr \eval \val \Then \EnvLoc_2,\EnvStore_2 \\*
		\EnvLoc_2,\EnvStore_2 &\then \addr = \alloc(\EnvStore_2) \\*
	}{
		\EnvLoc_2,\EnvStore_2 \then \code{$x$ $\tau$ <- $\expr$}
		\Then \EnvLoc_2\env{x \mapsto \addr},\EnvStore_2\env{\addr \mapsto \val}
	}
\]
It is semantically identical to a variable declaration statement followed by a
different redirection statement.


\subparagraph{Zero override} At (and only at) declaration, the zero value of a
variable can be overridden
\[
	\oprule[aligned]{
		\EnvSym               &\then x \in \Id \semiq \tau : \Type \\*
		\EnvLoc,\EnvStore     &\then \EnvLoc(x) = \varnothing \\*
		\EnvLoc,\EnvStore     &\then \expr \eval \val \Then \EnvLoc_2,\EnvStore_2 \\*
		\EnvLoc_2,\EnvStore_2 &\then \addr = \alloc(\EnvStore_2) \\*
	}{
		\EnvLoc_2,\EnvStore_2 \then \code{$x$ $\tau$ default $\expr$}
		\Then \EnvLoc_2\env{x \mapsto \addr},\EnvStore\env{\addr \mapsto \val}
	}
\]

\subparagraph{Initialization vs Zero override} Initialization and zero
override may seem identical, but they do two different things:
\begin{itemize}
	\item initialization: initializes the variable values with the one provided
		(without overriding its zero value, which remains equal to its type's
		zero value).
	\item zero override: overrides the variable's zero value (to which the
		variable is initialized to).
\end{itemize}
Both can be performed in a single statement
\[
	\oprule[aligned]{
		\EnvSym               &\then x \in \Id \semiq \tau : \Type \\*
		\EnvLoc,\EnvStore     &\then \EnvLoc(x) = \varnothing \\*
		\EnvLoc,\EnvStore     &\then \expr_z \eval \val_z \Then \EnvLoc_2,\EnvStore_2 \\*
		\EnvLoc_2,\EnvStore_2 &\then \expr \eval \val \Then \EnvLoc_3,\EnvStore_3 \\*
		\EnvLoc_3,\EnvStore_3 &\then \addr_z = \alloc(\EnvStore_3) \\*
		\EnvLoc_3,\EnvStore_3 &\then \addr = \alloc(\EnvStore_3) \\*
	}{
		\EnvLoc_3,\EnvStore_3 \then \code{$x$ $\tau$ default $\expr_z$ <- $\expr$}
		\Then \EnvLoc_3\env{x \mapsto \addr},\EnvStore\env{\addr \mapsto \val,\addr_z \mapsto \val_z} \\*
		\hfill\implies \code{$x$} \eval \val \land \Zero(x) = \val_z
	}
\]

\paragraph{Constant declaration}

\[
	\oprule[aligned]{
		\EnvLoc,\EnvStore     &\then x \in \Id \semiq \EnvLoc(x) = \varnothing \\*
		\EnvLoc,\EnvStore     &\then \expr \eval \val \Then \EnvLoc_2,\EnvStore_2 \\*
		\EnvLoc_2,\EnvStore_2 &\then \EnvLoc_2(\expr) = \addr_\expr \semiq \EnvStore_2(\addr_\expr) = \val \\*
	}{
		\EnvLoc_2,\EnvStore_2\then \code{const $x$ <- $\expr$} \Then \EnvLoc_2\env{x \mapsto \addr_\expr}
	}
\]

\comment{TODO: rettificare in caso scelgo deep immutability invece che shallow
immutability}

\paragraph{Alias declaration}

\[
	\oprule[aligned]{
		\EnvLoc,\EnvStore     &\then x \in \Id \semiq \EnvLoc(x) = \varnothing \\*
		\EnvLoc,\EnvStore     &\then \expr \eval \val \Then \EnvLoc_2,\EnvStore_2 \\*
		\EnvLoc_2,\EnvStore_2 &\then \EnvLoc_2(\expr) = \addr_\expr \semiq \EnvStore_2(\addr_\expr) = \val \\*
	}{
		\EnvLoc_2,\EnvStore_2\then \code{alias $\expr$ as $x$} \Then \EnvLoc_2\env{x \mapsto \addr_\expr}
	}
\]

\clearpage

\subsection{Patterns}

See \See{ch:patterns} for more details.

\paragraph{Patterns sets} The following sets group patterns to make rules more
readable
\[
	\begin{split}
		\Pattern         &= \PatternTerminal \cup \PatternProduct \cup \PatternIterator\\*
		\PatternTerminal &= \set{\pattern \mid \pattern := \text{any unary pattern or placeholder pattern}}\\*
		\PatternProduct  &= \set{\pattern \mid \pattern := \text{any product pattern}}\\*
		\PatternIterator &= \set{\pattern \mid \pattern := \text{any iterator pattern}}\\*
	\end{split}
\]

\paragraph{The environment set} The $\Env$ will be used to denote a set of
environments.\\*
For example, $\Env_\EnvLoc$ will denote the set of location environments.

\comment{TODO: rinominare in Decomposition?}

\paragraph{Deconstruction function} Given how patterns can be nested, a
recursive function will be used to describe the deconstruction.
\[
	\begin{split}
		\decons& : (\Env_\EnvLoc,\Env_\EnvStore,\Val,\Pattern) \to (\Env_\EnvLoc,\Env_\EnvStore) \\*
		\decons&(\EnvLoc,\EnvStore,\val, \pattern) \mapsto \begin{cases}
			(\EnvLoc\env{},\EnvStore\env{})                                                              &\text{ if } \pattern = \code{\_} \\*
			(\EnvLoc\env{a\mapsto\addr},\EnvStore\env{\addr\mapsto\val})                                 &\text{ if } \pattern = \code{$a$} \\*
			\bigcup_{i\in[1,n]}{\decons(\EnvLoc,\EnvStore,\val_i,\pattern_i)}                            &\text{ if } \val:\Product \\*
																										 &\quad\land \pattern = \code{<$\pattern_1$, $\dotsc$, $\pattern_n$>} \\*
			\decons(\EnvLoc,\EnvStore,\val',\pattern')                                                   &\text{ if } \val:\List \\*
																										 &\quad\land \pattern = \code{[$\pattern'$]} \\*
			\decons(\EnvLoc,\EnvStore,\val_k,\pattern_k)\cup\decons(\EnvLoc,\EnvStore,\val_v,\pattern_v) &\text{ if } \val:\Map \\*
																										 &\quad\land \pattern = \code{[$\pattern_k$ => $\pattern_v$]} \\*
		\end{cases}
	\end{split}
\]

In other words, the function tells that:
\begin{itemize}
	\item for a placeholder pattern, the resulting environments will be the
		same.
	\item for a unary pattern, the resulting environments will be extended with
		a new mapping for the specified value.
	\item for product pattern, the resulting environments will be extended with
		the mappings produced from applying the deconstruction function to
		every field of the product value.
	\item for a list pattern, the resulting environments will be extended with
		the mappings produced from applying the deconstruction function to a
		list's entry value.
	\item for a map pattern, the resulting environments will be extended with
		the mappings produced from applying the deconstruction function to a
		map's entry key and value pair.
\end{itemize}

\clearpage

\subsubsection{Terminal patterns}

\paragraph{Placeholder pattern}
\[
	\oprule{
		\EnvLoc,\EnvStore \then \expr \in \Expr \Then \EnvLoc_2,\EnvStore_2
	}{
		\EnvLoc_2,\EnvStore_2 \then \code{$\expr$ -> \_} \Then \EnvLoc_2,\EnvStore_2
	}
\]

\vspace{\baselineskip}

More formally, using the deconstruction function
\[
	\oprule{
		\pattern = \code{\_} \\*
		\EnvLoc,\EnvStore \then \expr\eval\val \Then \EnvLoc_2,\EnvStore_2 \\*
	}{
		\EnvLoc_2,\EnvStore_2 \then \code{$\expr$ -> $\pattern$} \Then \decons(\EnvLoc_2,\EnvStore_2,\val,\pattern) \\*
	}
\]

\paragraph{Unary pattern} The semantic rules are the same of redirections and
variable declarations, because they're secretly the same thing.

\clearpage

\subsubsection{Product patterns}

\paragraph{Positional}

In case only unary patterns are used:
\[
	\oprule[aligned]{
		\EnvLoc,\EnvStore      &\then n \geq 0 \semiq \forall_{i\in[1,n]}{\val_i\in\Val} \semiq \val = \bigtimes_{i \in [1,n]}{\val_i} \\*
		\EnvLoc,\EnvStore      &\then m \leq n \semiq \forall_{j\in[1,m]}{\pattern_j\in\PatternTerminal\setminus\set{\code{\_}}} \semiq \pattern = \code{<$\pattern_1$, $\dotsc$, $\pattern_m$>} \\*
		\EnvLoc,\EnvStore      &\then \expr \eval \val \Then \EnvLoc_2,\EnvStore_2 \\*
		\EnvLoc_2,\EnvStore_2  &\then \forall_{j\in[1,m]}{\addr_j = \alloc(\EnvStore_2)} \\*
	}{
		\EnvLoc_2,\EnvStore_2 \then \code{$\expr$ -> $\pattern$} \Then \EnvLoc_2\env{\forall_{j\in[1,m]}{\pattern_j\mapsto\addr_j}},\EnvStore_2\env{\forall_{j\in[1,m]}{\addr_j\mapsto\val_j}}
	}
\]

\vspace{\baselineskip}

More formally, in case of nested patterns
\[
	\oprule[aligned]{
		\EnvLoc,\EnvStore &\then n \geq 0 \semiq \forall_{i\in[1,n]}{\val_i\in\Val} \semiq \val = \bigtimes_{i \in [1,n]}{\val_i} \\*
		\EnvLoc,\EnvStore &\then m \leq n \semiq \forall_{j\in[1,m]}{\pattern_j\in\PatternTerminal\cup\PatternProduct} \semiq \pattern = \code{<$\pattern_1$, $\dotsc$, $\pattern_m$>} \\*
		\EnvLoc,\EnvStore &\then \expr \eval \val \Then \EnvLoc_2,\EnvStore_2 \\*
	}{
		\EnvLoc_2,\EnvStore_2 \then \code{$\expr$ -> $\pattern$} \Then \decons(\EnvLoc_2,\EnvStore_2,\val,\pattern)
		% \EnvLoc_2\env{\forall_{j\in[1,m]}{\pattern_j\mapsto\addr_j}},\EnvStore_2\env{\forall_{j\in[1,m]}{\addr_j\mapsto\val_j}}
	}
\]

\paragraph{Nominal}

In case only unary patterns are used:
\[
	\oprule[aligned]{
		\EnvLoc,\EnvStore     &\then n \geq 0 \semiq \forall_{i\in[1,n]}{x_i\in\mathcal{N}\subseteq\Id, \val_i\in\Val} \semiq \val = \tuple{x_1:\val_1,\dotsc,x_n:\val_n} \\*
		\EnvLoc,\EnvStore     &\then m \leq n \semiq \forall_{j\in[1,m]}{y_j\in\mathcal{M}\subseteq\mathcal{N},\pattern_j\in\PatternTerminal\setminus\set{\code{\_}}} \\*
		\EnvLoc,\EnvStore     &\then \pattern = \code{<$y_1$:$\pattern_1$, $\dotsc$, $y_m$:$\pattern_m$>} \\*
		\EnvLoc,\EnvStore     &\then \expr \eval \val \Then \EnvLoc_2,\EnvStore_2 \\*
		\EnvLoc_2,\EnvStore_2 &\then \forall_{j\in[1,m]}{\addr_j = \alloc(\EnvStore_2)} \\*
		                      &\EnvLoc_3   = \EnvLoc_2\env{\forall_{j\in[1,m]}{\pattern_j\mapsto\addr_j}}\\*
		                      &\EnvStore_3 = \EnvStore_2\env{\forall_{j\in[1,m]}{\exists_{x'\in\mathcal{N}=y_j\in\mathcal{M}}{\quad\addr_j\mapsto\EnvStore_2(\EnvLoc_2(x'))}}}\\*
	}{
		\EnvLoc_2,\EnvStore_2 \then \code{$\expr$ -> $\pattern$} \Then \EnvLoc_3,\EnvStore_3\\*
	}
\]

\nb{
	A notation abuse is used for $\EnvStore_2(\EnvLoc_2(x'))$, which describes
	the value in memory of the field named $x'$ of the record $\val$.
}

\vspace{\baselineskip}

More formally, in case of nested patterns
\[
	\oprule[aligned]{
		\EnvLoc,\EnvStore  &\then n \geq 0 \semiq \forall_{i\in[1,n]}{x_i\in\mathcal{N}\subseteq\Id, \val_i\in\Val} \semiq \val = \tuple{x_1:\val_1,\dotsc,x_n:\val_n} \\*
		\EnvLoc,\EnvStore  &\then m \leq n \semiq \forall_{j\in[1,m]}{y_j\in\mathcal{M}\subseteq\mathcal{N},\pattern_j\in\PatternTerminal\cup\PatternProduct} \\*
		\EnvLoc,\EnvStore  &\then \pattern = \code{<$y_1$:$\pattern_1$, $\dotsc$, $y_m$:$\pattern_m$>} \\*
		\EnvLoc,\EnvStore  &\then \expr \eval \val \Then \EnvLoc_2,\EnvStore_2 \\*
	}{
		\EnvLoc_2,\EnvStore_2 \then \code{$\expr$ -> $\pattern$} \Then \decons(\EnvLoc_2,\EnvStore_2,\val,\pattern)
	}
\]

\subparagraph{Implicitly named product pattern}

\[
	\oprule[aligned]{
		\EnvLoc,\EnvStore       &\then n \geq 0 \semiq \forall_{i\in[1,n]}{x_i\in\mathcal{N}\subseteq\Id, \val_i\in\Val} \semiq \val = \tuple{x_1:\val_1,\dotsc,x_n:\val_n} \\*
		\EnvLoc,\EnvStore       &\then m \leq n \semiq \forall_{j\in[1,m]}{\pattern_j\in\PatternTerminal\setminus\set{\code{\_}}} \\*
		\EnvLoc,\EnvStore       &\then \pattern = \code{<\_:$\pattern_1$, $\dotsc$, \_:$\pattern_m$>} \\*
		\EnvLoc,\EnvStore       &\then \expr \eval \val \Then \EnvLoc_2,\EnvStore_2 \\*
		\EnvLoc_2,\EnvStore_2   &\then \forall_{j\in[1,m]}{\addr_j = \alloc(\EnvStore_2)} \\*
		                        &\EnvLoc_3   = \EnvLoc_2\env{\forall_{j\in[1,m]}{\pattern_j\mapsto\addr_j}}\\*
		                        &\EnvStore_3 = \EnvStore_2\env{\forall_{j\in[1,m]}{\exists_{x'\in\mathcal{N}=\pattern_j\in\mathcal{M}}{\semiq\addr_j\mapsto\EnvStore_2(\EnvLoc_2(x'))}}}\\*
	}{
		\EnvLoc_2,\EnvStore_2 \then \code{$\expr$ -> $\pattern$} \Then \EnvLoc_3,\EnvStore_3\\*
	}
\]

\subsubsection{Iterator patterns}

\paragraph{List pattern}

\[
	\oprule[aligned]{
		\EnvLoc,\EnvStore &\then n \geq 0 \semiq \expr\eval\glist{\val_1,\dotsc,\val_n} \Then \EnvLoc_2,\EnvStore_2 \\*
						  &\forall_{i\in[1,n]}{(\EnvLoc_i,\EnvStore_i) = \decons(\EnvLoc_2,\EnvStore_2,\val_i,\pattern)}
	}{
		\EnvLoc_2,\EnvStore_2 \then \code{for [$\pattern$] in $\expr$ do} \Then \EnvLoc_i,\EnvStore_i\\*
	}
\]
At each iteration $i$, the $i$-th entry of the list gets deconstructed with
the pattern $\pattern$


\paragraph{Map pattern}

\[
	\oprule[aligned]{
		\EnvLoc,\EnvStore &\then n \geq 0 \semiq \expr\eval\map{\val_{k_1}\mapsto\val_{v_1},\dotsc,\val_{k_n}\mapsto\val_{v_n}} \Then \EnvLoc_2,\EnvStore_2 \\*
						  &\forall_{i\in[1,n]}{(\EnvLoc_i,\EnvStore_i) = \decons(\EnvLoc_2,\EnvStore_2,\val_{k_i},\pattern_k)} \cup \decons(\EnvLoc_2,\EnvStore_2,\val_{v_i},\pattern_v)
	}{
		\EnvLoc_2,\EnvStore_2 \then \code{for [$\pattern_k$ => $\pattern_v$] in $\expr$ do} \Then \EnvLoc_i,\EnvStore_i\\*
	}
\]
At each iteration $i$, the $i$-th key-value pair of the map gets deconstructed
with the patterns $\pattern_k$ and $\pattern_v$ respectively

\clearpage

\subsection{Expressions}

\subsubsection{Operators}

\[
	\oprule[aligned]{
		&\bop = \text{any binary operator} \\*
		\EnvLoc,\EnvStore     &\then \expr_1 \eval \val_1 \Then \EnvLoc_2,\EnvStore_2 \\*
		\EnvLoc_2,\EnvStore_2 &\then \expr_2 \eval \val_2 \Then \EnvLoc_3,\EnvStore_3 \\*
		\EnvLoc_3,\EnvStore_3 &\then \val_3 = \val_1 \bop \val_2 \\*
	}{
		\EnvLoc_3,\EnvStore_3\then \code{$\expr_1$ $\bop$ $\expr_2$} \eval \val_3 \Then \EnvStore_3\env{\alloc(\EnvStore_3) \mapsto \val_3}
	}
\]

\paragraph{Side-effects} Evaluating expressions may cause side effects (e.g.
a function invokation), hence new environments are declared.

\paragraph{Short-circuit evaluation} In some operations (e.g. boolean), some
expressions may be "short-circuited" (they won't be evaluated).\\*
For example: \code{false and someFunction()}, \texttt{someFunction()} will not
be evaluated.


\subsubsection{Invocation}

A function can be invoked using two approaches:
\begin{itemize}
	\item Positional \code{$f$($\expr_1$,$\dotsc$,$\expr_n$)}
	\item Nominal \code{$f$($x_1$:$\expr_1$,$\dotsc$,$x_n$:$\expr_n$)}
\end{itemize}

\paragraph{Positional}

\[
	\oprule[aligned]{
		\EnvSym               &\then n \geq 0 \semiq \forall_{i\in[1,n]}{x_i\in\Id,\tau_i:\Type} \semiq \tau_{n+1}\Type \\*
		\EnvSym               &\then \phi = (x_1:\tau_1,\dotsc,x_n:\tau_n)\to\tau_{n+1} \\*
		\EnvLoc,\EnvStore     &\then f\in\Expr \semiq f \eval \phi \Then \EnvLoc_2,\EnvStore_2 \\*
		\EnvLoc_2,\EnvStore_2 &\then \forall_{i\in[1,n]}{\expr_i\eval\val_i} \Then \EnvLoc_3,\EnvStore_3 \\*
		\EnvLoc_3,\EnvStore_3 &\then \forall_{i\in[1,n]}{\addr_i=\alloc(\EnvStore_3)} \\*
		                      &\EnvLoc_4 = \EnvLoc_3\env{\forall_{i\in[1,n]}{x_i\mapsto\addr_i}} \\*
		                      &\EnvStore_4 = \EnvStore_3\env{\forall_{i\in[1,n]}{\addr_i\mapsto\val_i}} \\*
	}{
		\EnvLoc_3,\EnvStore_3 \then \code{$f$($\expr_1$,$\dotsc$,$\expr_n$)} \eval \val_{\tau_{n+1}} \Then \EnvLoc_4,\EnvStore_4\env{\alloc(\EnvStore_4)\mapsto\val_{\tau_{n+1}}}
	}
\]
First, the expression $f$ to be invoked is evaluated (any expression that
evaluates to a function can be invoked); then all actual parameters $\expr_i$
are evaluated and bound to the respective formal parameter $x_i$.

\paragraph{Nominal}

\[
	\oprule[aligned]{
		\EnvSym               &\then n \geq 0 \semiq \forall_{i\in[1,n]}{x_i\in\mathcal{N}\subseteq\Id,\tau_i:\Type} \semiq \tau_{n+1}:\Type \\*
		\EnvSym               &\then \phi = (x_1:\tau_1,\dotsc,x_n:\tau_n)\to\tau_{n+1} \\*
		\EnvLoc,\EnvStore     &\then f\in\Expr \semiq f \eval \phi \Then \EnvLoc_2,\EnvStore_2 \\*
		\EnvLoc_2,\EnvStore_2 &\then \forall_{i\in[1,n]}{y_i\in\mathcal{N},\expr_i\eval\val_i} \Then \EnvLoc_3,\EnvStore_3 \\*
		\EnvLoc_3,\EnvStore_3 &\then \forall_{i\in[1,n]}{\addr_i=\alloc(\EnvStore_3)} \\*
		                      &\EnvLoc_4 = \EnvLoc_3\env{\forall_{i\in[1,n]}{\exists_{j\in[1,n]}{y_j=x_i \quad x_i\mapsto\addr_i}}} \\*
		                      &\EnvStore_4 = \EnvStore_3\env{\forall_{i\in[1,n]}{\exists_{j\in[1,n]}{y_j=x_i \quad \addr_i\mapsto\val_i}}} \\*
	}{
		\EnvLoc_3,\EnvStore_3 \then \code{$f$($y_1$:$\expr_1$,$\dotsc$,$y_m$:$\expr_m$)} \eval \val_{\tau_{n+1}} \Then \EnvLoc_4,\EnvStore_4\env{\alloc(\EnvStore_4)\mapsto\val_{\tau_{n+1}}}
	}
\]
First, the expression $f$ to be invoked is evaluated (any expression that
evaluates to a function can be invoked); then all actual parameters
$y_i:\expr_i$ are evaluated and bound to the formal parameter whose identifier
$x_i=y_j$, therefore the order of the parameter doesn't matter.

\subparagraph{Invocation labels} A function can declare an "invocation label"
for one or more formal parameters; the invocation label, if defined, has then
to be used instead of the formal parameter name when invoking the function.\\*
The rules are the same for a nominal invocation, except that, if a formal
parameter $x_i$ defines an invocation label $l_i$, then the actual parameter
name $y_j$ is bound to the formal parameter whose label is $l_i=y_j$.

\comment{TODO: se $\tau_1 = \Self$ o $\tau_{n+1} = \Self$}

\comment{TODO: opzionalità}

% \[
% 	\oprule[aligned]{
% 		\EnvSym               &\then n \geq 0 \semiq \forall_{i\in[1,n]}{x_i\in\mathcal{N}\subseteq\Id, \tau_i:\Type} \semiq \tau_{n+1}\Type \\*
% 		\EnvSym               &\then \phi = (x_1:\tau_1,\dotsc,x_n:\tau_n)\to\tau_{n+1} \\*
% 		\EnvLoc,\EnvStore     &\then f\in\Expr \semiq f \eval \phi \Then \EnvLoc_2,\EnvStore_2 \\*
% 		\EnvLoc_2,\EnvStore_2 &\then \forall_{i\in[1,n]}{\expr_i\eval\val_i} \Then \EnvLoc_3,\EnvStore_3 \\*
% 		\EnvLoc_3,\EnvStore_3 &\then m \leq n \semiq \forall_{j\in[1,m]}{y_j\in\mathcal{M}\subseteq\mathcal{N}} \\*
% 		\EnvLoc_3,\EnvStore_3 &\then \forall_{i\in[1,n]}{\addr_i=\alloc(\EnvStore_3)} \\*
% 		                      &\EnvLoc_4 = \EnvLoc_3\env{\forall_{y'\in\mathcal{M}}{y'\mapsto\alloc(\EnvStore_3)}} \cup \env{\forall_{x'\in\mathcal{N}\setminus\mathcal{M}}{x'\mapsto\alloc(\EnvStore_3)}} \\*
% 		                      &\EnvStore_4 = \EnvStore_3\env{\forall_{y'\in\mathcal{M}}{\EnvLoc(y')\mapsto\val_i}} \cup \env{\forall_{x'\in\mathcal{N}\setminus\mathcal{M}}{\EnvLoc(x')\mapsto\Zero(x')}}\\*
% 	}{
% 		\EnvLoc_3,\EnvStore_3 \then \code{$f$($y_1$:$\expr_1$,$\dotsc$,$y_m$:$\expr_m$)} \eval \val_{\tau_{n+1}} \Then \EnvLoc_4,\EnvStore_4\env{\alloc(\EnvStore_4)\mapsto\val_{\tau_{n+1}}}
% 	}
% \]

\chapter{Notation}

\section{Commonly used symbols}

\resizebox{0.55\linewidth}{!}{
\begin{tabular}{rll}
	\textbf{Symbol}                         & \textbf{Use}                     & \textbf{Example}s\\\hline
	$\bop$                                  & Generic binary operator          & $x \bop y$\\\hline
	$\narreq$                               & "Narrower than" operator         & $\Int \narreq \Real$ \\\hline
	$\sqcap$                                & "Least wide type" operator       & $\Int \sqcap \Real$ \\\hline
	$\then$                                 & Assumptions suffix               & $\EnvSym \then \val:\tau$ \\\hline
	$\Then$                                 & Environments modification prefix & $\expr \Then \EnvStore$\\\hline
	$\tau,\omega,\eta,\kappa,\nu$           & \multirow{2}{*}{A type}          & $\tau = \Int$\\
	or any lowercase greek letter           &                                  & $\eta=\Int \narreq \omega=\Real$\\\hline
	$\phi$                                  & Function type                    & $\phi=(\tau_1)\to\tau_2$\\\hline
	$\theta$                                & Tuple type                       & $\theta=\tau_1\times\tau_2$\\\hline
	$\rho$                                  & Record type                      & $\rho=\tuple{x:\tau_1,y:\tau_2}$\\\hline
	$\sigma$                                & Struct type                      & $\sigma=(\xi, \tuple{x:\tau_1,y:\tau_2})$\\\hline
	$\eval$                                 & "Evaluates to" operator          & $\expr\eval\val$\\\hline
	$\addr$                                 & Memory address                   & $\addr = 69$\\\hline
	$\Id$                                   & Identifiers set                  & \\\hline
	$a,b,x,y$                               & \multirow{2}{*}{Identifier}      & \\
	or any lower case latin letter          &                                  & \\\hline
	$\Expr$                                 & Expressions set                  & \\\hline
	$\expr$                                 & Expression                       & $\expr=\code{10 + 59}$\\\hline
	$\Val$                                  & Values set                       & \\\hline
	$\val$                                  & Value                            & $\val=69$\\\hline
	$\Zero$                                 & Zero set/function                & \\\hline
	$\Pattern$                              & Patterns set                     & \\\hline
	$\pattern$                              & Pattern                          & \\\hline
	$\PatternTerminal$                      & Terminal patterns set            & \\\hline
	$\PatternProduct$                       & Product patterns set             & \\\hline
	$\PatternIterator$                      & Iterator patterns set            & \\\hline
	$\Env$                                  & Environments set                 & \\\hline
	$\EnvSym$                               & General type environment         & \\\hline
	$\EnvVar$                               & Variables type environment       & \\\hline
	$\EnvConst$                             & Constants type environment       & \\\hline
	$\EnvFunc$                              & Functions type environment       & \\\hline
	$\EnvStucts$                            & Structs type environment         & \\\hline
	% $\EnvModules$                         & Modules type environment         & \\\hline
	$\EnvLoc$                               & Locations environment            & \\\hline
	$\EnvStore$                             & Store environment (memory)       & \\\hline
\end{tabular}
}

\section{Notations}

\resizebox{0.9\linewidth}{!}{
\begin{tabular}{rl}
	\multirow{2}{*}{\textbf{Notation}}                                                          & \textbf{Use} \\
	                                                                                            & \textbf{Example} \\\hline
	\multirow{2}{*}{$\code{monospace}$}                                                         & Verbatim code \\
	                                                                                            & \code{59 + 10} \\\hline
	\multirow{2}{*}{$\val : \tau$}                                                              & "$\val$ is a term of type $\tau$"\\
	                                                                                            & \code{69} : \Int \\\hline
	\multirow{2}{*}{$\expr \eval \val$}                                                         & "$\expr$ is evaluated to be equal to $\val$"\\
	                                                                                            & $\code{59 + 10} \eval 69$ \\\hline
	\multirow{2}{*}{$T[\tau:=\sigma]$}                                                          & "$T$ is a existential type with a type parameter $\tau$ of some type $\sigma$"\\
	                                                                                            & $\Map[\kappa:=\tau_\kappa,\nu:=\tau_\nu]$ \\\hline
	\multirow{2}{*}{$\exists\sigma.T$}                                                          & Concrete type $T$ with type parameter equal to $\sigma$\\
	                                                                                            & $\exists\Text.\exists\Int.\Map$ \\\hline
	\multirow{2}{*}{$\glist{\expr_1,\dotsc,\expr_n}$}                                           & A list of $n$ elements \\
	                                                                                            & $\glist{1,2,3}$ \\\hline
	\multirow{2}{*}{$\map{\expr_{k_1}\mapsto\expr_{v_1},\dotsc,\expr_{k_n}\mapsto\expr_{v_n}}$} & A Map of $n$ key-value pairs \\
	                                                                                            & $\map{a\mapsto1,b\mapsto2,c\mapsto3}$ \\\hline
	$(\tau_1,\dotsc,\tau_n)\to\tau_{n+1}$                                                       & Function type with $n$ formal parameters $\tau_i$ and return type $\tau_{n+1}$\\
	$(x_1:\tau_1,\dotsc,x_n:\tau_n)\to\tau_{n+1}$                                               & $(\Int,\Real)\to\Real$ \\\hline
	\multirow{2}{*}{$\tau_1 \times \dotsb \times \tau_n$}                                       & A Tuple type with $n$ fields of type $\tau_i$\\
	                                                                                            & $\Real \times \Real \times \Text$ \\\hline
	\multirow{2}{*}{$\tuple{x_1:\tau_1,\dotsc,x_n:\tau_n}$}                                     & A record type with $n$ fields with name $x_i$ and type $\tau_i$\\
	                                                                                            & $\tuple{\text{x}:\Real, \text{y}:\Real, \text{label}:\Text}$\\\hline
	\multirow{2}{*}{$(\xi, \tuple{x_1:\tau_1,\dotsc,x_n:\tau_n})$}                              & A struct type with typename $\xi$, and $n$ fields with name $x_i$ and type $\tau_i$\\
	                                                                                            & $(\text{Point}, \tuple{\text{x}:\Real, \text{y}:\Real, \text{label}:\Text})$\\\hline
	\multirow{2}{*}{$\env{x \mapsto y}$}                                                        & An environment in which $x$ maps to $y$\\
	                                                                                            & $\env{\addr \mapsto 69}$\\\hline
	\multirow{2}{*}{$E\env{}$}                                                                  & Environment "extension" (shorthand for $E\cup\env{}$)\\
	                                                                                            & $\EnvStore\env{\addr \mapsto 69}$\\\hline
\end{tabular}
}

\clearpage

\[
	\begin{split}
		A        \alpha   & \quad \text{alpha}   \\
		B        \beta    & \quad \text{beta}    \\
		\Gamma   \gamma   & \quad \text{gamma}   \\
		\Delta   \delta   & \quad \text{delta}   \\
		E        \epsilon & \quad \text{epsilon} \\
		Z        \zeta    & \quad \text{zeta}    \\
		H        \eta     & \quad \text{eta}     \\
		\Theta   \theta   & \quad \text{theta}   \\
		I        \iota    & \quad \text{iota}    \\
		K        \kappa   & \quad \text{kappa}   \\
		\Lambda  \lambda  & \quad \text{lambda}  \\
		M        \mu      & \quad \text{mu}      \\
		N        \nu      & \quad \text{nu}      \\
		\Xi      \xi      & \quad \text{xi}      \\
		O        o        & \quad \text{o}       \\
		\Pi      \pi      & \quad \text{pi}      \\
		P        \rho     & \quad \text{rho}     \\
		\Sigma   \sigma   & \quad \text{sigma}   \\
		T        \tau     & \quad \text{tau}     \\
		\Upsilon \upsilon & \quad \text{upsilon} \\
		\Phi     \phi     & \quad \text{phi}     \\
		X        \chi     & \quad \text{chi}     \\
		\Psi     \psi     & \quad \text{psi}     \\
		\Omega   \omega   & \quad \text{omega}   \\
	\end{split}
\]

\end{document}


